% Original section for integration into the main paper.
% Bibliography keys required: IshiiKernel, IshiiLie, BaumardSchneps,
% Rubin1991, GotoIshii2026, SanoVenjakob2026, MilneANT.
\section{The CM elliptic question after Ishii}
\label{sec:cm-audit}

This section identifies the arithmetic and Lie-theoretic statements
needed for the once-punctured CM elliptic curve, and proves the elementary
calculations at its two smallest split primes.
Let \(K\) be an imaginary quadratic field of class number one, let
\(\mathcal O_K\subset\operatorname{End}_K(E)\), and put \(X=E-\{O\}\).
For an ideal \(\mathfrak m\), write \(K(\mathfrak m)\) for the ray class
field of modulus \(\mathfrak m\), with no archimedean modulus. Fix a prime
\(p\geq5\) splitting as \((p)=\mathfrak p\bar{\mathfrak p}\).
Put \(K_\infty=\bigcup_nK(p^n)\), and let \(\Omega\) be the maximal
pro-\(p\) extension of \(K(p)\) unramified outside \(p\).
Ishii's Theorem A, p.~2, equivalently Theorem 2.14, p.~11,
identifies the fixed field of the outer kernel with \(K(E[p])\Omega\)
under three hypotheses: \(p\nmid h(K(p))\), exactly two primes of
\(K(p^2)\) above \(p\), and his Conjecture 2.3
\cite{Ishii}. The third hypothesis is a statement of freeness
for a particular arithmetic graded Lie algebra.

\subsection{The first Pollack relation and the arithmetic morphism}

Ishii's subsequent paper proves generation under the vanishing of
\(H^2_{\mathrm{\acute et}}(\mathcal O_K[1/p],\mathbb Q_p(m_1,m_2))\)
for every admissible pair \((m_1,m_2)\)
\cite[Theorem 1.2, p.~3]{IshiiLie}. Here
\(\mathbb Q_p(m_1,m_2)\) has character
\(\chi_{\mathfrak p}^{m_1}\chi_{\bar{\mathfrak p}}^{m_2}\), with
\(\chi_{\mathfrak p}\) and \(\chi_{\bar{\mathfrak p}}\) the two
characters of the CM Tate module, and admissibility means
\[
 m_1,m_2>0,\qquad m_1\equiv m_2\pmod {w_K},\qquad
 (m_1,m_2)\ne(1,1),
 \quad w_K=|\mathcal O_K^\times|.
\]
The same paper records that the earlier freeness assertion fails
when \(w_K=2\), and retains the two fields \(w_K=4,6\) as open cases
\cite[Former Conjecture 1.8, p.~4, and Remark 1.9, p.~5]{IshiiLie}.
We give the explicit derivation calculation behind this correction.
The relation is Pollack's; its derivations and their relation with
mixed elliptic motives are discussed by Baumard and Schneps
\cite[pp.~2--3]{BaumardSchneps}.

Let \(L=\mathbb Q_p\langle a,b\rangle_{\mathrm{Lie}}\) be the free
Lie algebra on \(a,b\), realized inside the free associative algebra
\(\mathbb Q_p\langle a,b\rangle\), with bracket \([u,v]=uv-vu\).
Set \(A_j=\operatorname{ad}_a^j(b)\). For even \(k\geq2\), define
a derivation by
\[
 \epsilon_k(a)=A_k,\qquad
 \epsilon_k(b)=\sum_{j=0}^{k/2-1}(-1)^j[A_j,A_{k-1-j}].
\]
A derivation is determined by its values on \(a,b\). Its bidegree
is the increase in the numbers of occurrences of \(a,b\); thus
\(\epsilon_k\) has bidegree \((k-1,1)\).

\begin{proposition}[The first Pollack relation]\label{prop:pollack-explicit}
These derivations satisfy
\[
 [\epsilon_4,\epsilon_{10}]-3[\epsilon_6,\epsilon_8]=0.
\]
Both brackets displayed here are nonzero.
\end{proposition}
\begin{proof}
The derivation identity and
\([a,[A_i,A_j]]=[A_{i+1},A_j]+[A_i,A_{j+1}]\)
give a telescoping sum
\(\epsilon_k([a,b])=0\).
For \(n\geq1\), repeated use of the derivation identity gives
\begin{align*}
 \epsilon_k(A_n)
 ={}&\sum_{r=0}^{n-1}\binom{n}{r+1}
              [A_{k+r},A_{n-1-r}]\\
 &+\sum_{j=0}^{k/2-1}(-1)^j
      \sum_{r=0}^{n}\binom nr
              [A_{j+r},A_{k-1-j+n-r}].
\end{align*}
Indeed, differentiating each occurrence of
\(\operatorname{ad}_a\) gives
\(\sum_{i=0}^{n-1}\operatorname{ad}_a^i
[A_k,A_{n-1-i}]\). Expanding this sum by the binomial theorem
and using
\(\sum_{i=r}^{n-1}\binom ir=\binom n{r+1}\)
gives the first line. Differentiating \(b\) gives the second.

Apply this formula to the two brackets and collect the coefficients
of \([A_j,A_{13-j}]\), \(0\leq j\leq6\), using
\([A_i,A_j]=-[A_j,A_i]\). The complete table is
\[
\begin{array}{c|rrrrrrr}
j&0&1&2&3&4&5&6\\ \hline
[\epsilon_4,\epsilon_{10}](a)&0&0&-6&-33&-75&-90&-42\\
{}[\epsilon_6,\epsilon_8](a)&0&0&-2&-11&-25&-30&-14 .
\end{array}
\]
Thus the proposed difference \(D\) annihilates \(a\).
It also annihilates \([a,b]\), since every \(\epsilon_k\) does.
Consequently \([a,D(b)]=0\).

For completeness, a finite associative polynomial commuting with
\(a\) contains only powers of \(a\). After removing those powers,
choose a nonzero word having the least number of initial \(a\)'s.
In \(Pa\) the corresponding word followed by \(a\) has that same
number of initial \(a\)'s; in \(aP\) every word containing \(b\)
has strictly more. This contradicts \(aP=Pa\).
Now \(D(b)\) has three occurrences of \(b\), so \(D(b)=0\).
The associative word \(a^{11}ba^2b\) has coefficient \(6\) in
\([\epsilon_4,\epsilon_{10}](a)\), and coefficient \(2\) in the
second bracket; hence neither bracket is zero.
\end{proof}

We specify how the relation reaches the arithmetic Lie algebra.
Let \(F^nG_K\) be the kernel of the outer action on the quotient
of the geometric pro-\(p\) fundamental group by its \((n+1)\)-st
lower central subgroup, and put
\(\mathfrak g_X=\bigoplus_{n>0}F^nG_K/F^{n+1}G_K\).
Ishii identifies \(\mathfrak g_X\otimes\mathbb Q_p\) with the
graded Lie algebra of the unipotent Galois image
\cite[Lemma 3.6, p.~14]{IshiiLie}. This gives an injection into
the outer derivations of \(L\) preserving the boundary class.
The boundary is \(z=[a,b]\), after the symplectic choice of free
generators. Ishii's Proposition 3.3(4), pp.~12--14, gives the
specific quotient map
\[
 \operatorname{Der}^{*}(L)/\mathbb Q_p\operatorname{ad}_z
       \ \xrightarrow{\ \sim\ }\ \operatorname{OutDer}^{*}(L).
\]
Here positive-degree derivations in \(\operatorname{Der}^{*}(L)\)
annihilate \(z\); the displayed kernel has bidegree \((1,1)\).
Therefore the quotient map is an isomorphism on each of the
bidegrees \((3,1),(5,1),(7,1),(9,1),(12,2)\) used below.

In bidegree \((m,1)\), a derivation annihilating \(z\) is a scalar
multiple of \(\epsilon_{m+1}\). Its value on \(a\) lies in the
one-dimensional Lie space spanned by \(A_{m+1}\): every bracket
with just one \(b\) is, by induction on its length, a multiple of
this iterated bracket. Subtract the corresponding multiple of
\(\epsilon_{m+1}\). The resulting derivation kills \(a\) and \(z\),
so it kills \(b\) by the centralizer argument above. This proves
the asserted one-dimensionality without choosing an arithmetic
normalization.

When \(w_K=2\), all four bidegrees \((3,1),(5,1),(7,1),(9,1)\)
are admissible. Write the images of four proposed nonzero generators
as \(c_m\epsilon_{m+1}\), with \(c_m\in\mathbb Q_p^\times\).
The relation then becomes
\[
 c_5c_7[\sigma_{(3,1)},\sigma_{(9,1)}]
 -3c_3c_9[\sigma_{(5,1)},\sigma_{(7,1)}]=0.
\]
The two commutators are linearly independent in the free Lie algebra
on the proposed generators: in its associative realization they
involve four distinct words. The injective arithmetic derivation
map therefore contradicts the proposed freeness. A vanishing
generator would already contradict that proposal.
For \(K=\mathbb Q(i)\) and \(\mathbb Q(\sqrt{-3})\), the index
\((3,1)\) is excluded by \(m_1\equiv m_2\pmod {w_K}\).
The calculation above gives no freeness proof in these two cases.

\subsection{Complex conjugation in the two-variable module}

Let \(\Gamma=\operatorname{Gal}(K_\infty/K(p))\simeq\mathbb Z_p^2\).
Complex conjugation interchanges its two local unit directions.
Choose generators \(\gamma_1,\gamma_2\) with
\(c\gamma_1c^{-1}=\gamma_2\), and identify
\[
 R=\mathbb Z_p[[\Gamma]]
   =\mathbb Z_p[[S,T]],\qquad S=\gamma_1-1,\quad T=\gamma_2-1 .
\]
For a module with the arithmetic conjugation action, the identity
is \(c(rm)=c(r)c(m)\), where \(c(S)=T\) and \(c(T)=S\).
The following calculation specifies the coefficient ring on which
the fixed subgroup is a module.

\begin{proposition}\label{prop:cm-plus}
Let \(M\) be a finitely generated compact \(R\)-module with a
continuous semilinear involution \(c\). Write \(M^+=\ker(c-1)\).
For \(m\in M^+\),
\[
 Sm\in M^+\quad\Longleftrightarrow\quad(S-T)m=0 .
\]
Consequently \(M^+\) is an \(R\)-submodule exactly when
\((S-T)M^+=0\). It is always a module over
\[
 R^c=\mathbb Z_p[[x,y^2]],\qquad
 x=(S+T)/2,\quad y=(S-T)/2 .
\]
When \(M^+\) is finite, \((S-T)M[1/p]=0\).
\end{proposition}
\begin{proof}
For \(cm=m\), semilinearity gives \(c(Sm)=Tm\). This proves
the first equivalence and the necessity of the second assertion.
Under the stated annihilation, multiplication by both \(S\) and
\(T\) preserves \(M^+\). It follows first for polynomials and then
by compactness and continuity for every element of \(R\).
Since \(p\) is odd, the change of variables \(S=x+y,T=x-y\) is
invertible over \(\mathbb Z_p\). The involution fixes \(x\) and
negates \(y\). A series is fixed exactly when its coefficients of
odd powers of \(y\) vanish. This proves the formula for \(R^c\),
and semilinearity proves its action on \(M^+\).

The two continuous projectors \((1+c)/2\) and \((1-c)/2\) give
the direct sum of the fixed and anti-fixed \(\mathbb Z_p\)-modules.
After inverting \(p\), finiteness kills the first summand.
Thus \(c=-1\) on \(M[1/p]\). Applying semilinearity to \(Sm\)
now gives \(-Sm=-Tm\), which proves the last assertion.
\end{proof}

This proposition identifies a precise arithmetic issue in formulating
the desired finiteness hypothesis. For example, \(R\) with its swapping
involution has \(1\in R^c\) but \(S\notin R^c\).
The module
\[
 M=\bigl(R/(S-T)\bigr)^{\oplus2},\qquad c(m)=-m,
\]
has \(M^+=0\); its action is semilinear because \(S=T\) in the
quotient. Nevertheless \(M[1/p]\) requires two generators at
the maximal ideal \((S,T)\) of \(R[1/p]\), since the quotient by
that ideal has dimension two over \(\mathbb Q_p\).
These are explicit module computations; they establish no occurrence
of either example as an arithmetic class-group module.

A second calculation isolates the information omitted by
characteristic ideals. Put
\[
 M_1=R/(S),\qquad
 M_2=R/(S)\oplus R/(S,T).
\]
Projection onto the first summand has kernel \(R/(S,T)\).
This kernel vanishes at every height-one prime of \(R\):
no such prime contains both \(S\) and \(T\), and either variable
annihilates the kernel. Thus \(M_1\) and \(M_2\) have the same
characteristic ideal \((S)\). After inverting \(p\), localization
at \((S,T)\) gives residue dimensions one and two, respectively;
the second module is noncyclic. Rubin's main conjecture computes
characteristic ideals of the relevant one-prime-ramified modules
in terms of elliptic units and Katz's \(p\)-adic \(L\)-functions
\cite{Rubin1991}. Applying that equality to the sought cyclicity
argument requires arithmetic control of the kernel illustrated
by this calculation.

\subsection{The full elliptic-unit radical for arbitrary split primes}
\label{subsec:cm-elliptic-radical}

This subsection requires only that \(K\) be an imaginary quadratic field
of class number one, that \(p\geq5\) split as
\((p)=\mathfrak p\bar{\mathfrak p}=(\pi)(\bar\pi)\), and that
\(E/K\) have complex multiplication by \(O_K\).
Put \(X=E-\{O\}\), \(L_n=K(p^n)\), and
\(L_\infty=\bigcup_{n\geq1}L_n\).
Let \(\Omega\) be the maximal pro-\(p\) extension of \(L_1\)
unramified outside \(p\), and set
\[
 H=\operatorname{Gal}(\Omega/L_\infty),\qquad A=H^{\mathrm{ab}}.
\]
All fields are taken inside the fixed algebraic closure.
Write \(\rho\) for the actual outer action of \(G_K\) on
\(\pi_1^{(p)}(X_{\bar K})\).
The proof of \cite[Lemma 2.13, p.~10]{Ishii} gives a field
\(\Omega^*\subseteq\Omega\) and the exact identity
\begin{equation}\label{eq:cm-radical-outer-field}
 \bar K^{\ker\rho}=K(E[p])\Omega^*,\qquad
 [K(E[p]):L_1]\ \text{prime to }p.
\end{equation}
The field \(L_\infty\) contains \(\mu_{p^\infty}\).
Also \(K(E[p^\infty])=K(E[p])L_\infty\), so
\(L_\infty\subseteq\Omega^*\) by intersecting with \(\Omega\).

Let \(C(L_n)\) be Rubin's elliptic-unit group, with the convention of
\cite[Definition 2.2, p.~8]{IshiiSoule}: it contains \(\mu(L_n)\)
and the appropriate Galois augmentation differences of theta values.
The convention using suitable twelfth roots gives the same
\(p\)-power radical because \(12\) is prime to \(p\).
Define the full radical field
\begin{equation}\label{eq:cm-full-radical}
 L_{\mathrm{ell}}=
 L_\infty\bigl(c^{1/p^r}:c\in C(L_n),\ n,r\geq1\bigr).
\end{equation}
It does not depend on choices of roots, since
\(\mu_{p^\infty}\subseteq L_\infty\).

\begin{proposition}\label{prop:cm-full-radical}
With the preceding definitions,
\[
             L_{\mathrm{ell}}\subseteq\Omega^*.
\]
Thus the actual outer kernel fixes every \(p\)-power root of every
Rubin elliptic unit in every finite ray layer \(L_n\).
This assertion holds for every split \(p\geq5\) and every such \(K,E\).
\end{proposition}

\begin{proof}
We give finite-level Kummer values throughout the argument.
Choose a \(p\)-compatible CM basis
\(\omega_{1,n}\in E[\mathfrak p^n]\),
\(\omega_{2,n}\in E[\bar{\mathfrak p}^n]\), and compatible
roots of unity \(\zeta_{p^r}\).
For a \(p\)-unit \(u\in L_\infty^\times\), all its \(p\)-power roots
belong to \(\Omega\).
Indeed, \(u\) belongs to a finite layer; after adjoining a sufficiently
large finite ray layer containing \(\mu_{p^r}\), its Kummer extension
has \(p\)-power degree and is unramified outside \(p\).
The normal closure over \(L_1\) is still a finite \(p\)-extension,
by taking all conjugate Kummer generators over a finite ray layer.
It is therefore contained in \(\Omega\).
For \(g\in H\), define \(k_u(g)\in\mathbb Z_p\) by
\[
 \frac{g(u^{1/p^r})}{u^{1/p^r}}
             =\zeta_{p^r}^{\,k_u(g)\bmod p^r}\quad(r\geq1).
\]
This value is independent of compatible root choices.
Moreover \(k_{uv}(g)=k_u(g)+k_v(g)\).
In particular \(k_u(g)=0\) means that \(g\) fixes every
\(p\)-power root of the specific finite-level element \(u\).

Put \(w=|O_K^\times|\) and
\[
 I=\{(m_1,m_2)\in\mathbb Z_{\geq1}^2-\{(1,1)\}:
                      m_1\equiv m_2\pmod w\}.
\]
The elliptic Soul\'e characters \(\kappa_{\boldsymbol m}\),
\(\boldsymbol m\in I\), factor through \(A\)
by \cite[Lemma 4.1, p.~16]{IshiiSoule}.
Their relation with \(\rho\) is given by the following explicit
composition, as in
\cite[Definition 2.5 and Theorem 2.6, p.~8]{Ishii}:
\[
 G_{K(E[p^\infty])}
 \xrightarrow{\rho}F^3\Gamma_{1,1}
 \xrightarrow[\sim]{\text{boundary-preserving lift}}
 F^3\Gamma_{1,1}^{\dagger}
 \longrightarrow\Psi_1^\dagger
 \xrightarrow[\sim]{H}
       \mathbb Z_p[[T_1,T_2]](1).
\]
The middle projection is induced by the maximal metabelian quotient
of the geometric fundamental group.
This composition is \(\alpha_{1,1}\), whose coefficient formula is
\begin{equation}\label{eq:cm-radical-geometric-moments}
 \alpha_{1,1}(g)=
 \sum_{\substack{r,s\geq0\\ r+s\geq2\ {\rm even}}}
 \frac{\kappa_{(r+1,s+1)}(g)}
      {(1-p^{r+s})\,r!\,s!}\,U_1^rU_2^s,\qquad
 U_i=\log(1+T_i).
\end{equation}
Consequently every element of \(\ker\rho\) has
\(\kappa_{\boldsymbol m}=0\) for every \(\boldsymbol m\in I\).
We now prove a stronger intermediate statement: every \(g\in H\)
with all these character values zero fixes \(L_{\mathrm{ell}}\).

Fix a nontrivial ideal \(\mathfrak a\) prime to \(6p\).
Let \(\theta_{\mathfrak a}\) be the rational function of
\cite[Definition 2.1, p.~7]{IshiiSoule}.
For nonzero \(p\)-power torsion \(P\), its value is a \(p\)-unit
in the ray field of the conductor of \(P\).
It is exactly invariant under \(O_K^\times=\operatorname{Aut}(E)\).
The distribution and reciprocity formulas used below are
\cite[Propositions 2.1--2.2, p.~7]{IshiiSoule}.
For
\[
             Y=\mathbb Z_p^2-p\mathbb Z_p^2,
\]
define a measure on each primitive residue cylinder by
\begin{equation}\label{eq:cm-cylinder-kummer}
 \mu_{\mathfrak a,g}\bigl((a,b)+p^n\mathbb Z_p^2\bigr)
   =k_{\theta_{\mathfrak a}(a\omega_{1,n}+b\omega_{2,n})}(g),
 \qquad p\nmid\gcd(a,b).
\end{equation}
These are exact \(\mathbb Z_p\)-values.
The \(p^2\) refinements at level \(n+1\) correspond to all
\([p]\)-preimages of \(a\omega_{1,n}+b\omega_{2,n}\).
The distribution identity
\[
  \prod_{\nu\in E[p]}\theta_{\mathfrak a}(Q+\nu)
                       =\theta_{\mathfrak a}(pQ)
\]
therefore proves finite additivity under refinement.
Every cylinder value lies in \(\mathbb Z_p\), so this defines a bounded
\(\mathbb Z_p\)-valued measure on \(Y\).

The exact finite-level identity
\cite[Lemma 3.1, p.~11]{IshiiSoule} is
\begin{equation}\label{eq:cm-primitive-moment-identity}
 \epsilon_{\boldsymbol m,n}^{
       N\mathfrak a-\chi^{\boldsymbol1-\boldsymbol m}
                                      (\sigma_{\mathfrak a})}
 =\prod_{\substack{0\leq a,b<p^n\\p\nmid\gcd(a,b)}}
 \theta_{\mathfrak a}(a\omega_{1,n}+b\omega_{2,n})^{
                            a^{m_1-1}b^{m_2-1}}
 \quad\text{in }L_\infty^\times/L_\infty^{\times p^n}.
\end{equation}
By definition the \(p^n\)-Kummer character of
\(\epsilon_{\boldsymbol m,n}\) is
\(\kappa_{\boldsymbol m}\bmod p^n\).
This formula applies also to \(m_1=1\) or \(m_2=1\), with
\(0^0=1\); thus it includes the two axes.
Since a polynomial is constant modulo \(p^n\) on the cylinders
of level \(n\), taking Kummer values in
\eqref{eq:cm-primitive-moment-identity} for all \(n\) gives
\begin{equation}\label{eq:cm-primitive-moment-value}
 \int_Y x^r y^s\,d\mu_{\mathfrak a,g}
 =\bigl(N\mathfrak a-\chi^{(-r,-s)}
                          (\sigma_{\mathfrak a})\bigr)
                  \kappa_{(r+1,s+1)}(g)
\end{equation}
whenever \(r,s\geq0\), \(r+s>0\), and \(r\equiv s\pmod w\).
These moments vanish by our choice of \(g\).

For \(\varepsilon\in O_K^\times\), the action on \(Y\) is
\[
 (x,y)\longmapsto
        \bigl(\iota_{\mathfrak p}(\varepsilon)x,
              \iota_{\bar{\mathfrak p}}(\varepsilon)y\bigr).
\]
Here the two unit eigenvalues are inverse roots of unity.
The exact invariance of \(\theta_{\mathfrak a}\) makes
\(\mu_{\mathfrak a,g}\) invariant under this action.
When \(r-s\) is not divisible by \(w\), averaging over this finite
unit group makes the moment of \(x^ry^s\) vanish.
Consequently all nonconstant monomial moments vanish.
The constant is recovered without deleting either axis.
For \(N_j=(p-1)p^j\), the function \(z^{N_j}\) converges uniformly
on \(\mathbb Z_p\) to the indicator of \(\mathbb Z_p^\times\).
Hence
\[
             x^{N_j}+y^{N_j}-x^{N_j}y^{N_j}
                         \longrightarrow1\quad\text{uniformly on }Y.
\]
All three integrals vanish, and so the constant moment vanishes.
Polynomial functions over \(\mathbb Q_p\) are uniformly dense in
continuous \(\mathbb Q_p\)-valued functions on \(Y\):
extend a continuous function by zero to the clopen complement in
\(\mathbb Z_p^2\) and use its two-variable Mahler expansion.
Boundedness now gives \(\mu_{\mathfrak a,g}=0\).
In view of \eqref{eq:cm-cylinder-kummer}, \(g\) fixes every \(p\)-power
root of every \(\theta_{\mathfrak a}(P)\) for
\((\mathfrak a,6p)=1\) and every nonzero \(p\)-power torsion point
\(P\).

We next treat exactly the additional auxiliary ideals allowed on
an axis.  Suppose \((\mathfrak a,6\mathfrak p)=1\), allowing
\(\bar{\mathfrak p}\mid\mathfrak a\).
Choose
\[
              c=1+6pN\mathfrak a,\qquad \mathfrak b=(c).
\]
Then \(\mathfrak b\) is prime to \(6p\mathfrak a\).
For a generator \(\alpha\) of \(\mathfrak a\),
\cite[Lemma 2.4, p.~7]{IshiiSoule} gives the rational-function
identity
\[
 \frac{\theta_{\mathfrak b}(P)^{N\mathfrak a}}
      {\theta_{\mathfrak b}(\alpha P)}
 =
 \frac{\theta_{\mathfrak a}(P)^{c^2}}
      {\theta_{\mathfrak a}(cP)}.
\]
If \(P\) has \(\mathfrak p\)-power order, all four terms are defined.
The preceding paragraph kills the Kummer values of the two
\(\theta_{\mathfrak b}\)-terms, including at \(\alpha P\), and gives
\begin{equation}\label{eq:cm-auxiliary-axis}
 c^2 k_{\theta_{\mathfrak a}(P)}(g)
                      =k_{\theta_{\mathfrak a}(cP)}(g).
\end{equation}

The points \(v_n=\bar\pi^n\omega_{1,n}\) satisfy
\(\pi v_{n+1}=v_n\).
Use them to define a bounded measure on \(\mathbb Z_p^\times\) by
\[
 \nu_{\mathfrak a,g}(x+p^n\mathbb Z_p)
                    =k_{\theta_{\mathfrak a}(xv_n)}(g).
\]
Its refinement identity is the \([\pi]\)-distribution formula;
\((\mathfrak a,\mathfrak p)=1\) is precisely its required
coprimality hypothesis.
Equation \eqref{eq:cm-auxiliary-axis} says
\(\nu_{\mathfrak a,g}(cB)=c^2\nu_{\mathfrak a,g}(B)\)
for every cylinder \(B\).
With convolution defined by
\([t]\nu(B)=\nu(t^{-1}B)\), the resulting equation is exactly
\begin{equation}\label{eq:cm-axis-injective-operator}
                    (c^2-[c^{-1}])\nu_{\mathfrak a,g}=0.
\end{equation}
The measure algebra has the decomposition
\[
 \mathbb Z_p[[\mathbb Z_p^\times]]
   =\prod_{\eta\in\widehat{\mu_{p-1}}}\mathbb Z_p[[T]].
\]
Since \(c\equiv1\pmod p\), in each factor the operator in
\eqref{eq:cm-axis-injective-operator} is
\(c^2-(1+T)^s\), where
\(s=\log(c^{-1})/\log(\gamma)\in\mathbb Z_p\) for a generator
\(\gamma\) of \(1+p\mathbb Z_p\).
Its constant coefficient is \(c^2-1\ne0\).
Each factor is an integral domain, so multiplication by this
operator is injective.
It follows that \(\nu_{\mathfrak a,g}=0\).
Thus \(g\) fixes all roots of all allowed theta values on the
\(\mathfrak p\)-axis.  The same proof with
\(\mathfrak p,\bar{\mathfrak p}\) interchanged handles the other
axis.

Finally, \cite[Proposition 2.9, pp.~9--10]{IshiiSoule} gives all
generators of \(C(L_n)\): roots of unity and the
\(\sigma-1\) differences, \(\sigma\in\operatorname{Gal}(L_n/K)\),
of
\[
 \begin{array}{c|c}
 \text{theta value}&\text{allowed auxiliary ideal}\\ \hline
 \theta_{\mathfrak a}(\omega_{1,n}+\omega_{2,n})
                       &(\mathfrak a,6p)=1\\
 \theta_{\mathfrak a}(\omega_{1,n})
                       &(\mathfrak a,6\mathfrak p)=1\\
 \theta_{\mathfrak a}(\omega_{2,n})
                       &(\mathfrak a,6\bar{\mathfrak p})=1 .
 \end{array}
\]
All conjugate values were included in the preceding cylinder
arguments.
Every root of unity in \(L_\infty\) has all its \(p\)-power roots
in \(L_\infty\): the \(p\)-primary part lies in
\(\mu_{p^\infty}\), and raising to \(p\) is invertible on the
prime-to-\(p\) part.
We have therefore proved
\begin{equation}\label{eq:cm-radical-kernel-inclusion}
 \bigcap_{\boldsymbol m\in I}\ker\kappa_{\boldsymbol m}
        \subseteq
 \ker\!\left(H\longrightarrow
                  \operatorname{Gal}(L_{\mathrm{ell}}/L_\infty)\right).
\end{equation}

In particular every element of the absolute outer kernel fixes
\(L_{\mathrm{ell}}\), by
\eqref{eq:cm-radical-geometric-moments}.
Thus \(L_{\mathrm{ell}}\subseteq K(E[p])\Omega^*\).
On the other hand \(L_{\mathrm{ell}}\subseteq\Omega\) by the
finite Kummer-extension argument at the start of the proof.
Since \(K(E[p])/L_1\) has prime-to-\(p\) degree while
\(\Omega/\Omega^*\) is pro-\(p\),
\[
                 \Omega\cap K(E[p])\Omega^*=\Omega^*.
\]
This exact intersection proves the desired inclusion.
\end{proof}

\begin{corollary}\label{cor:cm-radical-joint-kernel}
Put \(\bar A=\operatorname{Gal}(L_{\mathrm{ell}}/L_\infty)\),
with its canonical quotient map \(A\to\bar A\).
Then
\[
 \bigcap_{\boldsymbol m\in I}\ker
       (\kappa_{\boldsymbol m}:A\longrightarrow
                           \mathbb Z_p(\boldsymbol m))
                         =\ker(A\longrightarrow\bar A).
\]
In particular the geometric elliptic Soul\'e characters factor
through, and jointly separate, this full radical quotient.
\end{corollary}

\begin{proof}
One containment is \eqref{eq:cm-radical-kernel-inclusion}.
For the other, let \(g\in H\) fix \(L_{\mathrm{ell}}\).
For a fixed auxiliary ideal \((\mathfrak a,6p)=1\), the value
\(u_n=\theta_{\mathfrak a}(\omega_{1,n}+\omega_{2,n})\)
is a global unit in \(L_n\).
Writing \(d_n=[L_n:K]\), the identity
\[
 u_n^{d_n}
  =N_{L_n/K}(u_n)
       \prod_{\tau\in\operatorname{Gal}(L_n/K)}
                                (u_n^{\tau-1})^{-1}
                         \in C(L_n)
\]
holds because the norm is a global unit of the imaginary
quadratic field \(K\), hence a root of unity.
Thus \(d_n k_{u_n}(g)=0\) in the torsionfree group
\(\mathbb Z_p\), and \(k_{u_n}(g)=0\).
This applies to all conjugates of \(u_n\).

Let \(\kappa_{\boldsymbol m,\mathfrak a}\) denote the moment
character of the norm-compatible mixed-unit system \((u_n)_n\).
Its explicit reduction modulo \(p^n\) is the Kummer character of
\[
 \prod_{\tau\in\operatorname{Gal}(L_n/K)}
             \tau(u_n)^{\chi^{\boldsymbol m-\boldsymbol1}(\tau)}
\]
by \cite[Lemma 4.2, p.~17]{IshiiSoule}.
Consequently \(\kappa_{\boldsymbol m,\mathfrak a}(g)=0\).
The exact comparison
\cite[Proposition 4.3, pp.~17--18]{IshiiSoule} is
\begin{equation}\label{eq:cm-radical-reverse-scalar}
 w\left(\frac1{1-\iota^{\boldsymbol m-\boldsymbol1}(\pi)}
       +\frac1{1-\iota^{\boldsymbol m-\boldsymbol1}(\bar\pi)}
       -1\right)\kappa_{\boldsymbol m,\mathfrak a}
  =\bigl(N\mathfrak a-
           \chi^{\boldsymbol1-\boldsymbol m}(\sigma_{\mathfrak a})
    \bigr)\kappa_{\boldsymbol m}.
\end{equation}
Here
\(\iota^{(r,s)}(\alpha)=
 \iota_{\mathfrak p}(\alpha)^r
 \iota_{\bar{\mathfrak p}}(\alpha)^s\).
Both denominators are nonzero:
the corresponding algebraic numbers have complex absolute value
\(p^{(m_1+m_2-2)/2}>1\).
Choose \(\mathfrak a=(c)\) with a rational integer \(c>1\)
prime to \(6p\).
The ambiguity of the CM reciprocity action by an element of
\(O_K^\times\) disappears because
\(m_1-m_2\equiv0\pmod w\).
The scalar on the right of
\eqref{eq:cm-radical-reverse-scalar} is therefore exactly
\[
                 c^2-c^{\,2-m_1-m_2}\ne0.
\]
Equation \eqref{eq:cm-radical-reverse-scalar} now implies
\(\kappa_{\boldsymbol m}(g)=0\) for every
\(\boldsymbol m\in I\).
This proves the reverse containment.
\end{proof}

The proposition uses no assertion that an individual character is
nonzero or surjective.  Its conclusion identifies a specific
abelian quotient of \(H\) through an exact radical construction.
The further map from this quotient to a proposed presentation by
free Iwasawa summands is the additional module-theoretic step
discussed in \cite[Remark 4.9.1, p.~24]{Ishii}; no such
presentation is inferred here.

\subsection{Finite conjugation-fixed class modules and source-group freeness}
\label{subsec:cm-finite-plus}

Continue with \(K,p,L_n,L_\infty,\Omega,H,A\) from
Subsection~\ref{subsec:cm-elliptic-radical}.
Complex conjugation is denoted by \(c\).
Put
\[
 \mathcal X_2=\varprojlim_n\operatorname{Cl}(L_n)[p^\infty],
 \qquad
 \mathcal J_2=\varinjlim_n\operatorname{Cl}(L_n)[p^\infty],
\]
where the first transition maps are ideal norms and the second
are extension of ideals.
The finite-level class groups here are the full class groups.
Let \(\mathcal X_2^+=(1+c)\mathcal X_2/2\), considered as a compact
\(\mathbb Z_p\)-module.
This definition does not assume that the fixed subgroup is stable
under the two-variable Iwasawa algebra.
Write
\[
 G=\operatorname{Gal}(L_\infty/K),\quad
 \Gamma=\operatorname{Gal}(L_\infty/L_1)\cong\mathbb Z_p^2,\quad
 L_{\mathrm c}=L_1(\mu_{p^\infty}),\quad
 \Gamma_{\mathrm a}=\operatorname{Gal}(L_\infty/L_{\mathrm c})
                   \cong\mathbb Z_p .
\]
Ray reciprocity identifies
\[
 G=\bigl(\mathbb Z_p^\times\times\mathbb Z_p^\times\bigr)/O_K^\times,
 \qquad
 c(x,y)c^{-1}=(y,x),\qquad
 \chi_{\mathrm{cyc}}(x,y)=xy.
\]
The embedded \(O_K^\times\) has the form
\((\varepsilon,\varepsilon^{-1})\).
Thus \(G/\langle g(g^c)^{-1}:g\in G\rangle\) is precisely the
cyclotomic quotient \(\mathbb Z_p^\times\), and
\(\Gamma_{\mathrm a}\) is its pro-\(p\) kernel.

\begin{lemma}\label{lem:cm-full-class-kummer}
Let
\[
 \mathcal U_\infty=O_{L_\infty}^{\times}
       =\bigcup_n O_{L_n}^{\times},\qquad
 L_{\mathrm u}=L_\infty(u^{1/p^r}:
                         u\in\mathcal U_\infty,\ r\geq1).
\]
There is a \(G\)- and \(c\)-equivariant exact sequence
\begin{equation}\label{eq:cm-full-class-kummer}
 0\longrightarrow
 \mathcal U_\infty\otimes_{\mathbb Z}\mathbb Q_p/\mathbb Z_p
 \longrightarrow
 B:=\operatorname{Hom}_{\mathrm{cont}}(A,\mu_{p^\infty})
 \xrightarrow{\vartheta}\mathcal J_2
 \longrightarrow0 .
\end{equation}
Consequently the restriction map has kernel
\begin{equation}\label{eq:cm-unit-annihilator}
 N_{\mathrm u}:=
 \ker\!\left(A\longrightarrow
            \operatorname{Gal}(L_{\mathrm u}/L_\infty)\right)
 \cong\operatorname{Hom}(\mathcal J_2,\mu_{p^\infty}).
\end{equation}
\end{lemma}

\begin{proof}
We spell out \(\vartheta\); this also proves that the right-hand term
is the full direct-limit class group.
The same construction for a cyclotomic tower is
\cite[Proposition 3.3.9 and Corollary 3.3.14]{SharifiIwasawaNotes}.
Its valuation proof applies to the present tower because it
contains the cyclotomic tower.

Represent a finite-order character in \(B\) by a Kummer class
\(u\in L_\infty^\times/L_\infty^{\times p^r}\).
At every prime outside \(p\), its valuation is divisible by \(p^r\),
because the corresponding Kummer extension is unramified there.
Pass to a finite ray layer containing \(u\) and sufficiently high
in the cyclotomic direction that all valuations at primes over
\(p\) are divisible by \(p^r\) as well.
This is possible since their ramification indices grow without
bound.
At that layer there is a unique fractional ideal \(\mathfrak b\)
with
\[
                         (u)=\mathfrak b^{p^r}.
\]
Set \(\vartheta(u)=[\mathfrak b]\in\mathcal J_2\).
Replacing \(u\) by \(uv^{p^r}\) changes \(\mathfrak b\) by
the principal ideal \((v)\); going to a higher layer extends
\(\mathfrak b\).
Representing the same character at exponent \(p^{r+1}\) replaces
\(u\) by \(u^p\) and leaves \(\mathfrak b\) unchanged.
Thus the map is well defined and equivariant.

Every \(p^r\)-torsion ideal class at a finite layer has a
representative \(\mathfrak b\) with
\(\mathfrak b^{p^r}=(u)\); the associated Kummer extension is
unramified outside \(p\).
This proves surjectivity.
The class \([\mathfrak b]\) dies in the direct limit exactly when,
at a higher finite layer, \(\mathfrak b=(v)\).
Then \(uv^{-p^r}\) is a global unit.
This proves the asserted kernel.
The left Kummer map is injective because a \(p^r\)-th root of a
global unit, whenever it belongs to \(L_\infty\), is again a global
unit.
Taking the Pontryagin dual with values in \(\mu_{p^\infty}\)
gives \eqref{eq:cm-unit-annihilator}.
The Tate twist in this duality is retained: \(c\) acts by \(-1\)
on \(\mu_{p^\infty}\).
\end{proof}

\begin{lemma}\label{lem:cm-anticyclotomic-invariants}
Without a class-group finiteness assumption, one has
\[
 A^{\Gamma_{\mathrm a}}=0,\qquad
 H^2(H,\mathbb Q_p/\mathbb Z_p)=0 .
\]
\end{lemma}

\begin{proof}
Put \(H_{\mathrm c}=\operatorname{Gal}(\Omega/L_{\mathrm c})\).
These are the full maximal pro-\(p\) Galois groups unramified outside
\(p\) of their indicated infinite base fields.
To check the comparison, a finite Galois \(p\)-extension over either
infinite field descends to a finite subfield of that base,
which is a \(p\)-extension of \(L_1\).
After increasing that finite subfield, take the descended
Galois \(p\)-extension and adjoin all its conjugates over \(L_1\);
the resulting normal closure has \(p\)-power degree and is
unramified outside \(p\), and therefore belongs to \(\Omega\).
Thus we have the exact sequence of arithmetic groups
\begin{equation}\label{eq:cm-anticyclotomic-group-extension}
 1\longrightarrow H\longrightarrow H_{\mathrm c}
                 \longrightarrow\Gamma_{\mathrm a}
                 \longrightarrow1 .
\end{equation}

Iwasawa's theorem on weak Leopoldt for a cyclotomic
\(\mathbb Z_p\)-extension gives
\[
 H^2(H_{\mathrm c},\mathbb Q_p/\mathbb Z_p)=0;
\]
see \cite[Remark 4.3, author version p.~26]{VenjakobIwasawa},
which explicitly states the result for every number field and
refers to its cohomological formulation.
This theorem concerns the trivial coefficient
\(\mathbb Q_p/\mathbb Z_p\).
If it is formulated first for the full \(S\)-ramified Galois
group, its inflation to the maximal pro-\(p\) quotient in degree
two is injective.
Indeed, an invariant nonzero homomorphism from the kernel of
the maximal pro-\(p\) quotient to \(\mathbb F_p\) would produce
a larger pro-\(p\) quotient; the same assertion for
\(\mathbb Q_p/\mathbb Z_p\) follows by taking a nonzero
order-\(p\) multiple.
The five-term sequence proves the required injectivity.

Let \(M=H^1(H,\mathbb Q_p/\mathbb Z_p)=A^\vee\).
Since \(\operatorname{cd}_p(\Gamma_{\mathrm a})=1\),
Hochschild--Serre for
\eqref{eq:cm-anticyclotomic-group-extension} gives
\begin{equation}\label{eq:cm-anticyclotomic-HS}
 0\longrightarrow H^1(\Gamma_{\mathrm a},M)
 \longrightarrow H^2(H_{\mathrm c},\mathbb Q_p/\mathbb Z_p)
 \longrightarrow
 H^2(H,\mathbb Q_p/\mathbb Z_p)^{\Gamma_{\mathrm a}}
 \longrightarrow0 .
\end{equation}
For a generator \(\gamma\) of \(\Gamma_{\mathrm a}\),
\[
 H^1(\Gamma_{\mathrm a},M)=M/(\gamma-1)M,
 \qquad
 \bigl(M/(\gamma-1)M\bigr)^\vee=A^{\Gamma_{\mathrm a}}.
\]
The first vanishing follows.
Equation \eqref{eq:cm-anticyclotomic-HS} also makes the
\(\Gamma_{\mathrm a}\)-invariants of the indicated \(H^2(H,-)\)
zero.
A nonzero discrete \(p\)-primary module under a pro-\(p\) group
always has nonzero invariants: choose a nonzero element of order
\(p\), take the finite-dimensional span of its finite orbit,
and use the nonzero fixed subspace of a finite \(p\)-group
acting on a nonzero \(\mathbb F_p\)-space.
Thus \(H^2(H,\mathbb Q_p/\mathbb Z_p)=0\) as well.
No freeness of \(H\) was used.
\end{proof}

\begin{proposition}\label{prop:cm-finite-plus-free}
For the preceding CM ray tower, suppose that the explicitly
defined group \(\mathcal X_2^+\) is finite.
Then
\[
 \mathcal J_2=0,\qquad
 \operatorname{Gal}(L_{\mathrm u}/L_\infty)=A,
 \qquad H\text{ is a free pro-}p\text{ group}.
\]
In particular every finite-layer \(p\)-primary ideal class becomes principal
at some higher ray layer, and the maximal abelian subextension
of \(\Omega/L_\infty\) is generated by \(p\)-power roots of
global units.
\end{proposition}

\begin{proof}
First, every norm
\[
 \operatorname{Cl}(L_{n+1})[p^\infty]
           \longrightarrow\operatorname{Cl}(L_n)[p^\infty]
\]
is surjective.
Ray reciprocity identifies
\(\operatorname{Gal}(L_{n+1}/L_n)\) with the product of the
two principal-unit quotients, each of order \(p\).
The inertia at \(\mathfrak p\) contains the first factor and
the inertia at \(\bar{\mathfrak p}\) contains the second.
Hence \(L_{n+1}/L_n\) has no nontrivial everywhere-unramified
subextension.
By the Artin reciprocity description of the image of the ideal
norm, this is exactly the required surjectivity on class groups.
The same holds for all composite norm maps.
All maps commute with \(c\); averaging by \((1+c)/2\) proves
surjectivity on their fixed subgroups.

It follows that \(\mathcal X_2^+\) surjects onto every
\(\operatorname{Cl}(L_n)[p^\infty]^+\).
Since \(\mathcal X_2^+\) is finite, their orders are eventually
constant, and for all sufficiently large \(m\geq n\) the norm
\[
 N_{m/n}:\operatorname{Cl}(L_m)[p^\infty]^+
                       \xrightarrow{\sim}
                         \operatorname{Cl}(L_n)[p^\infty]^+
\]
is an isomorphism.
Let \(j_{m/n}\) be extension of ideals.
The exact identity
\[
                  N_{m/n}j_{m/n}=[L_m:L_n]
                               =p^{2(m-n)}
\]
holds on \(\operatorname{Cl}(L_n)\).
For fixed sufficiently large \(n\), this multiplication kills
the finite group \(\operatorname{Cl}(L_n)[p^\infty]^+\) when
\(m\) is sufficiently large.
Injectivity of \(N_{m/n}\) on the plus subgroup then gives
\(j_{m/n}=0\) there.
Consequently
\begin{equation}\label{eq:cm-direct-limit-plus-zero}
                         \mathcal J_2^+=0.
\end{equation}
This argument uses the exact norm and extension maps and does not
regard a plus subgroup as a two-variable Iwasawa submodule.

By \eqref{eq:cm-unit-annihilator} and the Tate twist in that
pairing, \eqref{eq:cm-direct-limit-plus-zero} says that \(c\)
acts as \(+1\) on \(N_{\mathrm u}\).
The subgroup \(N_{\mathrm u}\) is stable under \(G\), because
global units and their radical field are stable under \(G\).
For \(g\in G\) and \(a\in N_{\mathrm u}\), one therefore has
\[
                    ga=c(ga)=g^c\,ca=g^c a.
\]
Hence every \(g(g^c)^{-1}\) acts trivially on \(N_{\mathrm u}\).
The explicit ray-reciprocity quotient preceding the lemmas shows
that \(N_{\mathrm u}\subseteq A^{\Gamma_{\mathrm a}}\).
Lemma~\ref{lem:cm-anticyclotomic-invariants} makes this group zero.
Thus \(N_{\mathrm u}=0\).
The perfect Kummer duality
\eqref{eq:cm-unit-annihilator} now also yields
\(\mathcal J_2=0\), and proves the assertion about the maximal
abelian subextension.

It remains to prove freeness.
Equation \eqref{eq:cm-full-class-kummer} and
\(\mathcal J_2=0\) identify
\[
 H^1(H,\mu_{p^\infty})
       =\mathcal U_\infty\otimes_{\mathbb Z}
                                     \mathbb Q_p/\mathbb Z_p.
\]
The group on the right is \(p\)-divisible.
Since \(H\) acts trivially on \(\mu_{p^\infty}\), choosing
compatible roots of unity identifies it with
\(H^1(H,\mathbb Q_p/\mathbb Z_p)\).
The cohomology sequence for
\[
 0\longrightarrow\mathbb F_p
  \longrightarrow\mathbb Q_p/\mathbb Z_p
  \xrightarrow{p}\mathbb Q_p/\mathbb Z_p
  \longrightarrow0
\]
therefore gives \(H^2(H,\mathbb F_p)=0\), using the already
proved vanishing of \(H^2(H,\mathbb Q_p/\mathbb Z_p)\).
The cohomological freeness criterion for pro-\(p\) groups proves
that \(H\) is free.
For the same criterion in this arithmetic setting, see
\cite[Theorem 4.17 and its proof, author version p.~33]
{VenjakobIwasawa}.
\end{proof}

The group whose freeness has just been proved is the
source \(H=\operatorname{Gal}(\Omega/L_\infty)\).
To obtain a statement about an outer quotient, one must still
apply the specific restriction map defined by \(\rho\).
The radical calculation of the preceding subsection supplies
that map on the abelian quotient once the finite-level
elliptic-unit/global-unit index is compared.

\subsubsection{Character components and the exact Kummer twist}

Let
\[
 \Delta=\operatorname{Gal}(L_1/K)
       \cong\bigl((\mathbb F_p^\times)^2\bigr)/O_K^\times .
\]
We regard \(\Delta\) as the canonical prime-to-\(p\) factor of
\(G=\Delta\times\Gamma\).
All its characters take values in \(\mathbb Z_p^\times\), since
its exponent divides \(p-1\).
For such a character \(\chi\), write
\[
 e_\chi=\frac1{|\Delta|}
             \sum_{\delta\in\Delta}\chi(\delta)^{-1}\delta,
 \qquad
 \chi^c(\delta)=\chi(c\delta c^{-1}),\qquad
 \chi^\dagger=\omega_{\mathrm{cyc}}\chi^{-1},
\]
where \(\omega_{\mathrm{cyc}}=\chi_{\mathrm{cyc}}|_\Delta\).
For a \(c\)-stable set \(\mathcal O\) of characters, put
\[
 e_{\mathcal O}=\sum_{\chi\in\mathcal O}e_\chi,\qquad
 \mathcal O^\dagger=\{\chi^\dagger:\chi\in\mathcal O\}.
\]
The dagger operation commutes with \(c\).
In the pair notation
\(\chi_{r,s}(\delta_1,\delta_2)=\delta_1^r\delta_2^s\),
where \(r-s\equiv0\pmod w\), these operations are exactly
\[
                  (r,s)^c=(s,r),\qquad
                  (r,s)^\dagger=(1-r,1-s)
                                \pmod{p-1}.
\]

\begin{proposition}\label{prop:cm-component-unit-detection}
Let \(\mathcal O\) be a \(c\)-stable set of characters of
\(\Delta\), and suppose that
\((e_{\mathcal O}\mathcal X_2)^+\) is finite.
Then
\[
 e_{\mathcal O}\mathcal J_2=0,\qquad
 e_{\mathcal O^\dagger}N_{\mathrm u}=0.
\]
In particular, the canonical abelian field
\[
 L_{\mathcal O}
   =\Omega^{\,\ker(H\longrightarrow A
                       \xrightarrow{\,e_{\mathcal O^\dagger}\,}
                                e_{\mathcal O^\dagger}A)}
\]
satisfies \(L_{\mathcal O}\subseteq L_{\mathrm u}\).
\end{proposition}

\begin{proof}
The idempotent \(e_{\mathcal O}\) commutes with \(c\), with
extension of ideals, and with ideal norm.
The norm maps on
\((e_{\mathcal O}\operatorname{Cl}(L_n)[p^\infty])^+\)
are consequently surjective.
The eventual-isomorphism argument in the preceding proof,
using \(N_{m/n}j_{m/n}=p^{2(m-n)}\), gives
\[
                      (e_{\mathcal O}\mathcal J_2)^+=0.
\]

We record the component map in the Kummer pairing, including
its twist.  For \(a\in N_{\mathrm u}\),
\(b\in\mathcal J_2\), and \(\delta\in\Delta\), equivariance is
\[
              \langle\delta a,\delta b\rangle
                    =\delta\langle a,b\rangle
                    =\omega_{\mathrm{cyc}}(\delta)
                                         \langle a,b\rangle,
\]
where the pairing is written additively.
Thus a \(\psi\)-component of \(N_{\mathrm u}\) pairs with a
\(\chi\)-component of \(\mathcal J_2\) exactly when
\(\psi\chi=\omega_{\mathrm{cyc}}\).
The isomorphism \eqref{eq:cm-unit-annihilator} therefore
restricts to the exact isomorphism
\[
 e_{\mathcal O^\dagger}N_{\mathrm u}
       \cong
 \operatorname{Hom}(e_{\mathcal O}\mathcal J_2,
                                      \mu_{p^\infty}).
\]
Because the source on the right has only the \(c\)-minus part,
the group on the left is \(c\)-fixed.
It is stable under \(\Gamma\), since \(\Delta\) commutes with
\(\Gamma\).
For \(a\) in this subgroup and \(\gamma\in\Gamma\), one has
\(\gamma a=c(\gamma a)=\gamma^c a\).
The elements \(\gamma(\gamma^c)^{-1}\) generate
\(\Gamma_{\mathrm a}\), so this subgroup lies in
\(A^{\Gamma_{\mathrm a}}\).
Lemma~\ref{lem:cm-anticyclotomic-invariants} makes it zero.
Perfect duality then gives \(e_{\mathcal O}\mathcal J_2=0\).

Finally the restriction kernel
\(\ker(H\to\operatorname{Gal}(L_{\mathrm u}/L_\infty))\)
maps into \(N_{\mathrm u}\) in \(A\).
Its \(e_{\mathcal O^\dagger}\)-projection is zero, so it is
contained in the kernel defining \(L_{\mathcal O}\).
The Galois correspondence proves
\(L_{\mathcal O}\subseteq L_{\mathrm u}\).
\end{proof}

\subsection{The exact finite-level elliptic-unit comparison}
\label{subsec:cm-unit-index}

\begin{lemma}\label{lem:cm-unit-radical-comparison}
For every \(n\geq1\), the exact group \(C(L_n)\) in
\cite[Definition 2.2, p.~8]{IshiiSoule} has finite index in
\(O_{L_n}^{\times}\).
Consequently the two radical fields coincide:
\[
                         L_{\mathrm{ell}}=L_{\mathrm u}.
\]
\end{lemma}

\begin{proof}
We prove full rank with the defining theta values.
The formula below identifies their generators with the
Robert--Ramachandra invariants used in the logarithmic calculation.
Let \(F=L_n\), \(D=\operatorname{Gal}(F/K)\), and let
\(\chi\ne1\) be a complex character of \(D\).
Its conductor \(\mathfrak m\mid(p^n)\) is nontrivial, because
\(K\) has class number one.
Put \(F_{\mathfrak m}=K(\mathfrak m)\), and write
\(\chi_{\mathfrak m}\) for the induced primitive character.

Here is the required exact normalization.
For a lattice \(L\), the function
\[
 \varphi(z,L)=
       \bigl(e^{-\eta(z,L)z/2}\sigma(z,L)\bigr)^{12}\Delta(L)
\]
in \cite[(2), p.~255]{Oukhaba2009} is the fundamental theta
function in \cite[Definition 3.1(i), p.~6]{GotoIshii2026}.
The quotient formula in
\cite[Definition 3.1(ii), equation (1), p.~6]{GotoIshii2026}
gives, with precisely the rational function used by Ishii,
\[
 \theta_{\mathfrak a}(z;L)=
              \frac{\varphi(z,L)^{N\mathfrak a}}
                   {\varphi(z,\mathfrak a^{-1}L)}.
\]
Write \(L=\Omega_E O_K\), \(\mathfrak m=(\beta)\), and
\(\tau=\Omega_E/\beta\bmod L\).
Scaling both \(z\) and \(L\) leaves \(\varphi\) unchanged:
the factor \(\sigma^{12}\) scales by the twelfth power,
\(\Delta\) by its inverse, and \(\eta(z,L)z\) is invariant.
Thus \((\tau,L)\) may be replaced by \((1,\mathfrak m)\).

Let \(e_{\mathfrak m}>0\) generate
\(\mathfrak m\cap\mathbb Z\).
Oukhaba's invariants and reciprocity, equations (3)--(4) on
p.~255, state
\[
 \varphi_{\mathfrak m}(\sigma_{\mathfrak a})
      =\varphi(1,\mathfrak a^{-1}\mathfrak m)^{e_{\mathfrak m}},
 \qquad
 \varphi_{\mathfrak m}(1)^{\sigma_{\mathfrak a}}
      =\varphi_{\mathfrak m}(\sigma_{\mathfrak a}).
\]
Combining the displayed formulas gives the exact algebraic identity
\begin{equation}\label{eq:cm-exact-unit-convention}
 \theta_{\mathfrak a}(\tau)^{e_{\mathfrak m}}
       =\varphi_{\mathfrak m}(1)^{
                                  N\mathfrak a-\sigma_{\mathfrak a}}.
\end{equation}
The integer \(e_{\mathfrak m}\), which may be divisible by \(p\),
has been retained.

For \(u\in F^\times\), define
\[
 \ell_\chi^F(u)=
       \sum_{\sigma\in D}\chi(\sigma)^{-1}
                                    \log|\sigma(u)|^2.
\]
Then
\(\ell_\chi^F(u^\delta)=\chi(\delta)\ell_\chi^F(u)\).
For \(u\in F_{\mathfrak m}^\times\), one has
\[
 \ell_\chi^F(u)=[F:F_{\mathfrak m}]
                       \ell_{\chi_{\mathfrak m}}^{F_{\mathfrak m}}(u).
\]
The Kronecker limit formula
\cite[(6), p.~256]{Oukhaba2009}, applied at the primitive
conductor \(\mathfrak m\), gives
\[
 \ell_{\chi_{\mathfrak m}}^{F_{\mathfrak m}}
                         (\varphi_{\mathfrak m}(1))
   =-12r_{\mathfrak m}e_{\mathfrak m}
                          L'(0,\chi_{\mathfrak m}^{-1})\ne0 .
\]
Here \(r_{\mathfrak m}\) counts roots of unity of \(K\)
congruent to \(1\bmod\mathfrak m\).
The simple zero and nonzero derivative for a nontrivial primitive
character are also recorded explicitly in
\cite[\S2.3, p.~6, before (2.3)--(2.4)]{Oukhaba2003}.

Take \(c_0=1+6p\) and \(\mathfrak a=(c_0)\).
This auxiliary ideal is prime to \(6\mathfrak m\), and
\(N\mathfrak a-\chi_{\mathfrak m}(\sigma_{\mathfrak a})\ne0\):
the first term is the real integer \(c_0^2>1\) and the second
is a root of unity.
Equation \eqref{eq:cm-exact-unit-convention} therefore gives
\[
 \ell_\chi^F(\theta_{\mathfrak a}(\tau))
  =-12r_{\mathfrak m}[F:F_{\mathfrak m}]
    \bigl(N\mathfrak a-\chi_{\mathfrak m}(\sigma_{\mathfrak a})\bigr)
                         L'(0,\chi_{\mathfrak m}^{-1})\ne0.
\]
Choose \(\delta\in D\) with \(\chi(\delta)\ne1\).
The element
\[
                       v_\chi=\theta_{\mathfrak a}(\tau)^{\delta-1}
\]
is a generator allowed in Ishii's exact definition of \(C(F)\).
Indeed, \(F_{\mathfrak m}\subseteq F\), so the norm in that
definition is the identity.
The unit-injectivity requirement for \(\mathfrak m\) holds:
the norm of a difference of distinct roots of unity in an
imaginary quadratic field has prime divisors only \(2\) and \(3\),
while \(\mathfrak m\) is supported above \(p\geq5\).
For one-prime conductor the augmentation is still a global unit,
as in that definition; this also follows from the equal valuation
at all primes over that conductor in
\cite[(5), p.~256]{Oukhaba2009}.
Its \(\chi\)-logarithm is
\[
                \ell_\chi^F(v_\chi)
                   =(\chi(\delta)-1)
                      \ell_\chi^F(\theta_{\mathfrak a}(\tau))\ne0 .
\]

Dirichlet's logarithmic embedding of
\(O_F^\times/\mu(F)\) spans the real augmentation space of
\(\mathbb R[D]\).
This follows by choosing the embeddings extending the fixed
complex embedding of \(K\), one from each complex place.
After complexification, each nontrivial character of \(D\)
occurs exactly once.
The image of \(C(F)\) is \(D\)-stable, and the preceding
construction gives its nonzero projection to every such line.
It therefore has the full unit rank.
Since it is a subgroup of the finitely generated unit group
and already contains its torsion, its index is finite.

Finally, for any \(u\in O_{L_n}^\times\) there is an integer
\(d>0\) with \(u^d\in C(L_n)\).
If \(g\in H\) fixes \(L_{\mathrm{ell}}\), its exact Kummer
character satisfies \(d\,k_u(g)=0\) in \(\mathbb Z_p\).
Torsionfreeness gives \(k_u(g)=0\).
Thus \(g\) fixes all \(p\)-power roots of every global unit,
even when \(p\mid d\).
The reverse inclusion of radical fields follows from
\(C(L_n)\subseteq O_{L_n}^\times\), proving
\(L_{\mathrm{ell}}=L_{\mathrm u}\).
\end{proof}

\begin{theorem}\label{thm:cm-finite-plus-abelian}
Let \(\mathcal O\) be a \(c\)-stable set of characters of
\(\Delta=\operatorname{Gal}(K(p)/K)\), and suppose that
\((e_{\mathcal O}\mathcal X_2)^+\) is finite.
With
\(\mathcal O^\dagger=\{\omega_{\mathrm{cyc}}\chi^{-1}:
                                         \chi\in\mathcal O\}\),
the canonical field
\[
 L_{\mathcal O}=
 \Omega^{\,\ker(H\longrightarrow A
                \xrightarrow{\,e_{\mathcal O^\dagger}\,}
                              e_{\mathcal O^\dagger}A)}
\]
is contained in \(\Omega^*\).
Equivalently the image in \(A\) of the outer kernel has
zero projection to \(e_{\mathcal O^\dagger}A\).

For the set of all characters, finiteness of
\(\mathcal X_2^+\) gives
\[
 \Omega^{[H,H]}=L_{\mathrm{ell}}=L_{\mathrm u}
                   \subseteq\Omega^*,\qquad
 \bigcap_{\boldsymbol m\in I}\ker(\kappa_{\boldsymbol m}|_A)=0.
\]
The source \(H\) is then free pro-\(p\), by
Proposition~\ref{prop:cm-finite-plus-free}.
\end{theorem}

\begin{proof}
Proposition~\ref{prop:cm-component-unit-detection} gives
\(L_{\mathcal O}\subseteq L_{\mathrm u}\).
The preceding lemma identifies
\(L_{\mathrm u}=L_{\mathrm{ell}}\), and
Proposition~\ref{prop:cm-full-radical} places this field in
\(\Omega^*\).
These are actual inclusions inside \(\Omega\), so their
composition proves the component assertion.
For all characters, Proposition~\ref{prop:cm-finite-plus-free}
identifies \(L_{\mathrm u}\) with the maximal abelian
subextension \(\Omega^{[H,H]}\).
Corollary~\ref{cor:cm-radical-joint-kernel} then gives the
stated joint-character injectivity.
\end{proof}

\subsection{A source quotient detected by the CM depth-two measure}
\label{subsec:cm-measure-class-two}

Assume \(\mathcal X_2^+\) is finite and retain
\(H=\operatorname{Gal}(\Omega/L_\infty)\) and
\(A=H^{\mathrm{ab}}\).  The proved CM results make \(H\)
free pro-\(p\) and make the joint elliptic Soul\'e map
injective on \(A\).

There is a second useful measure realization of this map.
Put \(V=T_p(E)\simeq\mathbb Z_p^2\), written additively.
By \cite[Lemma 2.4, Definition 2.5 and Theorem 2.6,
pp.~7--8]{Ishii}, the canonical boundary-preserving
lift of the outer action gives
\[
 \alpha:A\longrightarrow
       \mathbb Z_p[[V]](1).
\]
Here \(\mathbb Z_p[[V]]\) denotes bounded measures on \(V\).
If \(U_i=\log(1+T_i)\), its exponential transform is
\[
 f_h(U_1,U_2)
 =\sum_{\substack{r,s\geq0\\r+s\geq2\ {\rm even}}}
 \frac{\kappa_{(r+1,s+1)}(h)}
      {(1-p^{r+s})r!s!}\,U_1^rU_2^s .
\]
The coefficients outside the admissible CM weight set
vanish.  The already proved joint-kernel equality makes
\(\alpha\) injective.

Define two automorphisms of \(V\times V\) by
\[
 t_1(x,y)=(x,x+y),\qquad t_2(x,y)=(x+y,y),
\]
and define the bounded-measure operator
\[
 \mathcal J=(t_1)_*+(t_2)_*-\operatorname{id}.
\]
The one-dimensional restriction of this operator occurs in
\cite[equations (29) and (32), p.~233]{LuqueNovelliThibon}.
To specify the comparison, send their coefficient
\(\phi_{n+1}\) to \(U^n/n!\), so that
\(\Phi(t)\) becomes \(\exp(tU)\), the exponential transform
of the point measure \(\delta_t\). Send an ordinary bracket
\([f,g]\) to \(f(U)g(V)-g(U)f(V)\). Their equation (32)
then becomes
\[
 \delta_s\wedge\delta_{s+t}
 +\delta_{s+t}\wedge\delta_t-\delta_s\wedge\delta_t
                 =\mathcal J(\delta_s\wedge\delta_t).
\]
For every \(v\in V\), the map \(i_v(t)=tv\) preserves
addition and gives
\((i_v\times i_v)_*\mathcal J_1
 =\mathcal J(i_v\times i_v)_*\).
The source calculation below proves the displacement formula
for the CM arithmetic action on the two-dimensional Tate module.

Use the alternating tensor inclusion to define
\[
 J_A:\widehat{\bigwedge}^2 A
 \longrightarrow\mathbb Z_p[[V\times V]](2),\qquad
 a\wedge b\longmapsto
 \mathcal J\bigl(\alpha(a)\otimes\alpha(b)
                         -\alpha(b)\otimes\alpha(a)\bigr).
\]
Let \(N_2^{\rm CM}\) be the inverse image of
\(\ker J_A\) under the canonical map
\(\gamma_2H\to\gamma_2H/\gamma_3H
 \simeq\widehat{\bigwedge}^2A\).

\begin{proposition}\label{prop:cm-decomposable-commutators}
For \(a,b\in A\) whose images in \(A[1/p]\) are linearly
independent, \(J_A(a\wedge b)\ne0\). Thus every such pair
has a nonzero commutator in \(Q_2^{\rm CM}\).
\end{proposition}
\begin{proof}
Write \(f=\alpha(a)\), \(g=\alpha(b)\) in the exponential
coordinates above. Both have zero constant and linear terms.
The exponential transform is injective: the Amice transform
identifies bounded measures with integral power series, and
the substitution \(T_i=\exp(U_i)-1\) has the inverse
\(U_i=\log(1+T_i)\) over \(\mathbb Q_p\).
The transform of \(J_A(a\wedge b)\) is
\begin{align*}
 &f(u)g(u+v)-g(u)f(u+v)\\
 &\quad+f(u+v)g(v)-g(u+v)f(v)-f(u)g(v)+g(u)f(v).
\end{align*}
Its coefficient linear in \(v_i\) is
\(f(u)\partial_i g(u)-g(u)\partial_i f(u)\), since every
term containing \(f(v)\) or \(g(v)\) has order at least two.
If the displayed transform vanishes, both derivatives of
\(g/f\) vanish in \(\operatorname{Frac}(\mathbb Q_p[[U_1,U_2]])\).
Embed this field into \(\mathbb Q_p((U_1))((U_2))\).
Every nonzero power series is invertible there after extracting
its first nonzero \(U_2\)-coefficient; the embedding commutes
with formal derivatives. The \(U_2\)-derivative forces the
ratio to belong to \(\mathbb Q_p((U_1))\), and its
\(U_1\)-derivative forces it to belong to \(\mathbb Q_p\),
by coefficient comparison in characteristic zero. Thus
\(g=cf\). Injectivity of \(\alpha[1/p]\) gives \(b=ca\),
contradicting independence. The commutator assertion follows
from the quotient theorem below.
\end{proof}

\begin{theorem}\label{thm:cm-measure-class-two}
The subgroup \(N_2^{\rm CM}\) is closed and stable under
the arithmetic action, and contains the outer
kernel.  The quotient
\[
                       Q_2^{\rm CM}=H/N_2^{\rm CM}
\]
is torsion-free of class at most two, has abelianization
\(A\), and has central commutator module
\(\operatorname{im}J_A\).
Its class is exactly two for every class-number-one
CM field in the stated range.
\end{theorem}

\begin{proof}
Let \(\Pi\) be the geometric pro-\(p\) fundamental group
of \(E-\{O\}\), with the generators \(x_1,x_2\)
of \cite[Lemma 2.1]{Ishii}. Put \(C=[\Pi,\Pi]\), \(C_r=\gamma_rC\), and
\(M=C/C_2=\mathbb Z_p[[V]]\bar z\), with
\(z=[x_2,x_1]\) as in \cite[Lemma 2.1]{Ishii}.
The closed subgroup \(C\) of the free pro-\(p\) group
\(\Pi\) is free pro-\(p\).  The integral class-two Hall
basis, first on finite subsets of a convergent free
basis and then by inverse limits, gives
\[
 C_2/C_3=\widehat{\bigwedge}^{\,2}M,
 \qquad \gamma_2H/\gamma_3H=
            \widehat{\bigwedge}^{\,2}A.
\]
The exterior square identifies with the antisymmetric
submodule of the completed tensor square by
\(a\wedge b\mapsto a\otimes b-b\otimes a\).
Since \(p\) is odd, half the exterior projection is
its inverse on this submodule.  Thus
\(C_2/C_3\) is a closed submodule of
\(\mathbb Z_p[[V\times V]](2)\).

For \(v\in V\), write \(e_v=v\bar z\), and let
\(T_v\) denote translation on \(M\) and simultaneous
translation on its tensor square.  Let \(f=s(h)\)
be the canonical special lift of \(h\in H\), and set
\(a=\alpha(h)\bar z\).  This lift fixes \(z\) exactly:
its exponent on the inertia generator is its
abelian determinant, which is one.  For \(q\in\Pi\)
with abelian image \(v\), put
\(k_f(q)=f(q)q^{-1}\in C\).
The crossed-homomorphism identity and Ishii's defining
formula on \(x_1,x_2\) give
\[
                 k_f(q)\bmod C_2=(T_v-1)a.
\]
This follows first on words and then by continuity.
Because \(C\) is normally generated by \(z\), the identity
\[
 f(qzq^{-1})=k_f(q)(qzq^{-1})k_f(q)^{-1}
\]
shows that \(f\) acts trivially on \(M\).  Its first
nontrivial action on \(C\) is the continuous map
\[
 \delta_a:M\longrightarrow C_2/C_3,
 \qquad \delta_a(e_v)=[T_va-a,e_v].
\]
In particular, it depends only on \(\alpha(h)\), and
\begin{equation}\label{eq:cm-delta-translation}
 \delta_a(T_vb)=T_v\delta_a(b)+[T_va-a,T_vb].
\end{equation}
The formula follows on point measures by substitution
and extends to all bounded measures by continuity.

Take another lift \(g=s(k)\), with
\(b=\alpha(k)\bar z\).  Their metabelian actions
commute.  In the class-two group \(C/C_3\), the BCH
law involves only the invertible denominator \(2\).
Use
\(k_{fg}(q)=f(k_g(q))k_f(q)\) and its reversed-order
counterpart.  Their unknown central displacements
cancel, and the surviving terms are
\begin{align*}
 k_{[f,g]}(q)\bmod C_3
  ={}&\delta_a((T_v-1)b)-\delta_b((T_v-1)a)\\
     &+[(T_v-1)b,(T_v-1)a].
\end{align*}
Indeed a special automorphism with all displacements
in \(C_2\) acts trivially on \(C/C_3\), by the
conjugate-of-\(z\) identity.  Hence replacing \(q\)
by \(gf(q)\) when passing from \(fg,gf\) to
\([f,g]\) does not change this displayed displacement.
Writing \(A_0=T_va,B_0=T_vb\), the bracket terms
introduced by \eqref{eq:cm-delta-translation} simplify as
\[
 [A_0-a,B_0]-[B_0-b,A_0]+[B_0-b,A_0-a]
                         =[A_0,B_0]-[a,b].
\]
We therefore obtain the exact identity
\[
 k_{[f,g]}(q)\bmod C_3
       =(T_v-1)(\delta_ab-\delta_ba+[a,b]).
\]
From the point-measure formula for \(\delta\),
\begin{align*}
 \delta_ab-\delta_ba+[a,b]
 &=\iint\bigl([e_x,e_{x+y}]+[e_{x+y},e_y]
                         -[e_x,e_y]\bigr)\,da(x)\,db(y)\\
 &=\mathcal J(a\wedge b).
\end{align*}
This proves the integral group comparison
\begin{equation}\label{eq:cm-integral-J-comparison}
 k_{[s(h),s(k)]}(q)\bmod C_3
      =(T_v-1)J_A(h^{\rm ab}\wedge k^{\rm ab}).
\end{equation}

For completeness, let \(\mathcal D_r\) be the special
automorphisms acting trivially on \(V\) whose
displacements all lie in \(C_r\).
An element of \(\mathcal D_1\) acts trivially on
\(C_2/C_3\), whereas an element of \(\mathcal D_2\)
acts trivially on \(C/C_3\).  Comparing their product
displacements and using \([C,C_2]\subset C_3\) proves
\([\mathcal D_1,\mathcal D_2]\subset\mathcal D_3\).
Consequently \(s(\gamma_3H)\subset\mathcal D_3\),
and additivity of displacement on \(\mathcal D_2\)
extends \eqref{eq:cm-integral-J-comparison} continuously
to every \(h\in\gamma_2H\):
\[
 k_{s(h)}(q)\bmod C_3=(T_v-1)J_A(h\bmod\gamma_3H).
\]
For \(q=x_1\), the operator \(T_v-1\) is multiplication
by
\[
 (1+T_1)(1+S_1)-1=T_1+S_1+T_1S_1
\]
in the power-series domain
\(\mathbb Z_p[[T_1,T_2,S_1,S_2]]\), so it is injective.
An outer-kernel element belongs to \(\gamma_2H\) by
abelian detection and has \(s(h)=1\), since the
canonical special lift is inverse to the outer
projection in \(F^3\).  The last identity therefore
places it in \(N_2^{\rm CM}\).

Here is also the exact arithmetic compatibility.
An arithmetic outer automorphism has a representative
preserving \(\langle z\rangle\); any two such
representatives differ by a power of
\(\operatorname{inn}(z)\), which commutes with every
\(s(h)\).  Uniqueness in \(F^3\) makes the special lift
equivariant.  On \(M\), conjugation is the linear
Tate action with determinant twist.  Both shears
commute with simultaneous linear changes on \(V\),
so \(J_A\) is equivariant with the displayed second
twist.  Thus its inverse-image kernel is stable
under the arithmetic action.  It is closed, and its
image is compact.  The resulting exact central extension
\[
 1\longrightarrow\operatorname{im}J_A
 \longrightarrow Q_2^{\rm CM}\longrightarrow A
 \longrightarrow1
\]
has commutator subgroup exactly \(\operatorname{im}J_A\).
Both end modules are \(p\)-torsion-free, so the group
is torsion-free and has the claimed abelianization.

It remains to prove that the commutator module is nonzero.
Let \(I_K\) denote the admissible weights, and put
\[
 \mathcal S=\{\boldsymbol m\in I_K:
  1\leq m_i\leq4,\ (5,5)-\boldsymbol m\in I_K\}.
\]
These are exactly the moments that can contribute
to the bidegree-\((5,5)\) bracket coefficient.
The ray-class description and two CM Tate characters
in \cite[pp.~4--5]{Ishii} show that
\(\Gamma=\operatorname{Gal}(L_\infty/L_1)\) has open
rank-two principal-unit image. Write
\(\chi_{\boldsymbol m}=\chi_1^{m_1}\chi_2^{m_2}\)
for the corresponding product of the two Tate characters. Choose \(\gamma\)
so that \(\chi_{\boldsymbol m}(\gamma)\), for
\(\boldsymbol m\in\mathcal S\), are pairwise distinct.
This is possible because after the \(p\)-adic logarithm
the finitely many equality loci are proper linear
subspaces of an open \(\mathbb Z_p^2\).
For \(\boldsymbol m_0=(2,2),(3,3)\), define the integral
finite-difference operator
\[
 P_{\boldsymbol m_0}(X)=
 \prod_{\substack{\boldsymbol m\in\mathcal S\\
                   \boldsymbol m\ne\boldsymbol m_0}}
       (X-\chi_{\boldsymbol m}(\gamma)).
\]
Both \(\kappa_{22}\) and \(\kappa_{33}\) are nonzero
by \cite[Theorem 2.8(2) and the following Soul\'e-finiteness
paragraph, p.~9]{Ishii}.
Choose \(a_0,b_0\in A\) on which these respective
characters are nonzero, and set
\(a=P_{22}(\gamma)a_0,b=P_{33}(\gamma)b_0\).
Equivariance gives
\[
 \kappa_{\boldsymbol m}(P(\gamma)d)
   =P(\chi_{\boldsymbol m}(\gamma))\kappa_{\boldsymbol m}(d).
\]
Thus all unwanted moments in \(\mathcal S\) vanish,
while the two desired moments remain nonzero.

In exponential moment variables \(U=(x,y),V=(z,t)\),
\(\mathcal JF=F(U+V,V)+F(U,U+V)-F(U,V)\).
Direct multiplication gives
\[
 J(XY,X^2Y^2)
       =(xy-zt)(xt+zt+yz)(xt+xy+yz)\ne0.
\]
The bidegree-\((3,3)\) polynomial component of
\(J_A(a\wedge b)\), corresponding to bracket weight
\((5,5)\), is this polynomial multiplied by
\[
 \frac{\kappa_{22}(a)\kappa_{33}(b)}
      {(1-p^2)(1-p^4)(2!)^2}\ne0.
\]
The definition of \(\mathcal S\) excludes every other
contribution.  Hence \(\operatorname{im}J_A\ne0\),
proving exact class two for every \(K\) in the statement.
\end{proof}

\subsection{Centralizers and free class-two source subgroups}
\label{subsec:cm-class-two-subgroups}

Retain the hypotheses and notation of
Theorem~\ref{thm:cm-measure-class-two}, and write
\[
 Q=Q_2^{\rm CM},\qquad Z=\operatorname{im}J_A,
 \qquad \pi:Q\longrightarrow A.
\]
The following statements use the proved integral map from the
actual source group to this quotient.

\begin{proposition}\label{prop:cm-class-two-centralizers}
The centre of \(Q\) is exactly \(Z\).  For every \(q\in Q\)
with \(a=\pi(q)\ne0\),
\[
 C_Q(q)=\pi^{-1}\bigl(A\cap\mathbb Q_p a\bigr).
\]
The intersection is taken inside \(A[1/p]\), is a free
\(\mathbb Z_p\)-module of rank one, and the displayed
centralizer is, noncanonically, \(Z\times\mathbb Z_p\).
In particular, every abelian subgroup of \(Q\) has image
of \(\mathbb Z_p\)-rank at most one in \(A\).
\end{proposition}

\begin{proof}
The commutator of two elements with abelian images \(a,b\)
is \(J_A(a\wedge b)\).  Proposition~\ref{prop:cm-decomposable-commutators}, together
with injectivity of \(\alpha[1/p]\), gives the exact criterion
\[
 J_A(a\wedge b)=0\quad\Longleftrightarrow\quad
 b\in\mathbb Q_pa\qquad(a\ne0).
\]
This proves the asserted centralizer formula.

Here is the integral topology of the line intersection.
Choose a continuous \(\mathbb Z_p\)-linear functional
\(\ell:A\to\mathbb Z_p\) with \(\ell(a)\ne0\).
Such a functional exists because the compact torsion-free
abelian pro-\(p\) group \(A\) is a product of copies of
\(\mathbb Z_p\).  If \(\lambda a\in A\), then
\(\lambda\ell(a)\in\mathbb Z_p\), which bounds the
denominator of \(\lambda\).  The allowed scalars form a
nonzero closed \(\mathbb Z_p\)-submodule of this bounded
rank-one lattice.  They are consequently \(p^e\mathbb Z_p\)
for some integer \(e\).  Thus \(A\cap\mathbb Q_pa\)
has a generator \(a_0\).  Choose \(q_0\in Q\) above
\(a_0\).  The continuous powers of \(q_0\) identify
its closed procyclic group with \(\mathbb Z_p\), and
\((z,t)\mapsto zq_0^t\) is an isomorphism from
\(Z\times\mathbb Z_p\) onto the displayed centralizer.

The module \(A[1/p]\) has dimension at least two: for
example, the moment-isolated elements constructed below
have distinct nonzero initial homogeneous terms.
Every nonzero \(a\) therefore has a \(\mathbb Q_p\)-independent
\(b\), so an element above \(a\) is not central.
Hence the centre is precisely \(Z\).  The assertion about
abelian subgroups follows from the same centralizer formula.
\end{proof}

\begin{corollary}\label{cor:cm-three-generator-class-two}
Let \(r\le3\), and let \(h_1,\ldots,h_r\in H\) have
\(\mathbb Q_p\)-linearly independent images in \(A[1/p]\).
For the closed subgroup \(U\) they generate, the map
\(U\to Q\) has kernel exactly \(\gamma_3U\), and
\(U/\gamma_3U\) is the free class-two pro-\(p\) group
on these \(r\) generators.
\end{corollary}

\begin{proof}
Let \(B\) be the \(\mathbb Z_p\)-span of their abelian
images.  Every bivector in the exterior square of a
vector space of dimension at most three over
\(\mathbb Q_p\) is decomposable.  The preceding
derivative argument thus makes
\(J_A:\bigwedge^2B\to Z\) injective.
Map the free pro-\(p\) group \(P_r\) to \(H\) using
the chosen elements.  The composite to \(Q\) kills
\(\gamma_3P_r\); its maps on the abelian quotient and
on \(\gamma_2P_r/\gamma_3P_r\simeq\bigwedge^2\mathbb Z_p^r\)
are injective.  Its kernel is therefore exactly
\(\gamma_3P_r\).  The image of that subgroup in \(U\)
is \(\gamma_3U\), which proves both conclusions.
\end{proof}

\begin{proposition}\label{prop:cm-countable-class-two}
There is a closed subgroup \(U\subset H\), generated
by a sequence tending to one, for which
\[
 U\cap N_2^{\rm CM}=\gamma_3U,
 \qquad U/\gamma_3U\simeq P/\gamma_3P,
\]
where \(P\) is the free pro-\(p\) group on a countably
infinite convergent sequence of generators.
Thus the complete free class-two quotient of this
actual source subgroup is detected.  Every finite
initial subsequence gives the corresponding conclusion
in finite rank.
\end{proposition}

\begin{proof}
For \(i\ge1\), put \(k_i=2^i\) and
\(d_i=2k_i-2\).  The character
\(\kappa_{(k_i,k_i)}\) is nonzero by
\cite[Theorem 2.8(2) and the Soul\'e-finiteness paragraph,
p.~9]{Ishii}; this is the ordinary \(k_i\)-th
Tate twist, with \(k_i\ge2\).

Let
\[
 \mathcal S_i=\{\boldsymbol m\in I:
                              |\boldsymbol m|\le2k_i\}.
\]
Put \(\Gamma_2=\operatorname{Gal}(L_\infty/L_1)\), and choose \(\gamma_i\in\Gamma_2\) for which the finitely
many characters \(\chi_{\boldsymbol m}(\gamma_i)\),
\(\boldsymbol m\in\mathcal S_i\), are pairwise distinct.
The open rank-two principal-unit image of the two CM
Tate characters, established by the ray-class morphism
in \cite[pp.~4--5]{Ishii}, permits this choice:
the forbidden loci after taking logarithms are a finite
union of proper linear subspaces.  Put
\[
 P_i(X)=\prod_{\substack{\boldsymbol m\in\mathcal S_i\\
                         \boldsymbol m\ne(k_i,k_i)}}
                 (X-\chi_{\boldsymbol m}(\gamma_i)).
\]
Choose \(b_i^0\in A\) with
\(\kappa_{(k_i,k_i)}(b_i^0)\ne0\), and set
\(b_i=P_i(\gamma_i)b_i^0\).  Equivariance shows that
all its moments in \(\mathcal S_i\), except the desired
one, vanish, and that the desired one remains nonzero.

Choose a lift \(\widetilde h_i\in H\) of \(b_i\), and
put \(h_i=\widetilde h_i^{p^i}\), with abelian image
\(a_i=p^ib_i\).  This sequence tends to one: in every
finite quotient of \(H\), a sufficiently large power
\(p^i\) kills every element.  The exponential transform
\(f_i\) of \(\alpha(a_i)\) has the form
\[
 f_i(U_1,U_2)=c_i(U_1U_2)^{k_i-1}
                    +\text{terms of total degree greater than }d_i,
 \qquad c_i\in\mathbb Q_p^\times.
\]
In particular, the \(d_i\) are strictly increasing.

Let \(P\) be free on the convergent sequence
\(x_i\to1\).  Its abelianization is
\(\prod_{i\ge1}\mathbb Z_p x_i\), and the integral
class-two Hall calculation gives
\[
 \gamma_2P/\gamma_3P
       =\prod_{i<j}\mathbb Z_p[x_i,x_j].
\]
These assertions follow by taking inverse limits of
the corresponding finite-generator, finite-quotient
Hall presentations.  The continuous map \(x_i\mapsto h_i\)
induces an injective abelian map: in a nonzero completed
linear combination of the \(f_i\), the least index
with nonzero coefficient supplies a unique nonzero
initial degree.

The initial homogeneous term of \(J(f_i,f_j)\) is
\[
 c_ic_j J\bigl((U_1U_2)^{k_i-1},
              (U_1U_2)^{k_j-1}\bigr),
\]
of degree
\[
 d_i+d_j=2^{i+1}+2^{j+1}-4.
\]
It is nonzero by the derivative computation above:
the two displayed monomials have nonconstant ratio.
These degrees are pairwise distinct for \(i<j\), by
uniqueness of binary expansion.  A nonzero completed
alternating combination therefore again has a unique
least initial degree, which cannot cancel.
For each total degree there are only finitely many
pairs to consider, so this argument applies to the
completed product, not only to finite combinations.
Consequently the composite
\(P\to H\to Q\) induces an injection on both its
abelian quotient and its class-two central quotient.
Its kernel is exactly \(\gamma_3P\).

Let \(U\) be the image of \(P\) in \(H\).  It is compact
and hence closed.  The image of \(\gamma_3P\) is
\(\gamma_3U\), by the commutator definition and
compactness.  The exact kernel just calculated gives
\(U\cap N_2^{\rm CM}=\gamma_3U\) and the asserted
isomorphism.  No injectivity assertion for the full
map \(P\to H\) is needed for these conclusions.
\end{proof}

The subgroups in Proposition~\ref{prop:cm-countable-class-two}
depend on the moment-isolating choices.  The surrounding
quotient \(Q_2^{\rm CM}\) remains the canonical arithmetic
quotient of Theorem~\ref{thm:cm-measure-class-two}.


\subsection{A mixed-unit measure and the exact lifting obstruction}
\label{subsec:cm-nonabelian-gate}

Keep the notation of Subsection~\ref{subsec:cm-elliptic-radical}, and put
\[
 N_{\mathrm{ell}}=\ker\!\left(A\longrightarrow
                  \operatorname{Gal}(L_{\mathrm{ell}}/L_\infty)\right),
 \qquad \overline A=A/N_{\mathrm{ell}},\qquad
 \Lambda_2=\mathbb Z_p[[\Gamma]].
\]
The following construction uses only the mixed torsion points whose
two coordinates are units.  The passage from its moments to the
full radical is justified by an explicit scalar formula, including
the two boundary families of weights.

\begin{lemma}\label{lem:cm-mixed-measure-embedding}
Fix an integer \(c_0>1\) prime to \(6p\), put
\(\mathfrak a=(c_0)\), and set
\[
 u_n=\theta_{\mathfrak a}(\omega_{1,n}+\omega_{2,n}).
\]
There is a continuous injective \(G\)-equivariant homomorphism
\[
 \overline\Phi_{c_0}:\overline A
            \lhook\joinrel\longrightarrow \mathbb Z_p[[G]](1).
\]
Here the target action is
\(\gamma\cdot\mu=\chi_{\mathrm{cyc}}(\gamma)[\gamma]\mu\);
translation is defined by
\([\gamma]\mu(B)=\mu(\gamma^{-1}B)\).
In particular \(\overline A\) is torsion-free over \(\Lambda_2\).
\end{lemma}

\begin{proof}
Let \(G_n=\operatorname{Gal}(L_n/K)\).  The exact norm relation
of \cite[Proposition 2.5, p.~8]{IshiiSoule} gives
\(N_{L_{n+1}/L_n}(u_{n+1})=u_n\).  Its root-of-unity factors are
one because no nonidentity unit of \(O_K\) is congruent to one
modulo \(p\), for \(p\geq5\).
For \(h\in H\), define a bounded measure by the finite-level
Kummer values
\begin{equation}\label{eq:cm-mixed-measure}
 \Phi_{c_0}(h)(\sigma\operatorname{Gal}(L_\infty/L_n))
                      =k_{\sigma(u_n)}(h)\qquad(\sigma\in G_n).
\end{equation}
The norm identity says exactly that the value of a cylinder is
the sum of the values of its refinements.  Thus these values give
an element of \(\mathbb Z_p[[G]]\).  Additivity of Kummer characters
makes \(\Phi_{c_0}\) a continuous homomorphism on \(A\).
For a lift of \(\gamma\in G\), the defining root equation gives
\[
 k_u(\gamma h\gamma^{-1})
   =\chi_{\mathrm{cyc}}(\gamma)k_{\gamma^{-1}u}(h).
\]
This proves the stated action and its independence of the lift,
since inner conjugation acts trivially on \(A\).

We identify its kernel using the moment formula for all weights.  By the proof of
Corollary~\ref{cor:cm-radical-joint-kernel}, a positive integral
power of each \(u_n\) belongs to \(C(L_n)\).  Consequently
\(N_{\mathrm{ell}}\subseteq\ker\Phi_{c_0}\).
Conversely, suppose \(\Phi_{c_0}(h)=0\).
For every \(\boldsymbol m=(m_1,m_2)\in I\), the exact moment
formula of \cite[Lemma 4.2, p.~17]{IshiiSoule} gives
\[
 \kappa_{\boldsymbol m,\mathfrak a}(h)
   =\int_G\chi^{\boldsymbol m-\boldsymbol1}(\sigma)\,
                                 d\Phi_{c_0}(h)(\sigma)=0.
\]
In this source \(\kappa_{\boldsymbol m,\mathfrak a}\) is defined
from this same mixed norm system for every \(\boldsymbol m\in I\).
When one of \(m_1,m_2\) equals one, the proof of
\cite[Proposition 4.3, pp.~17--18]{IshiiSoule} takes the norm
to the corresponding one-prime ray tower and includes its Euler
factor.  Its conclusion still reads
\begin{equation}\label{eq:cm-mixed-all-weight-scalar}
 w_K\,\frac{1-p^{r+s}}{(1-A_{\boldsymbol m})(1-B_{\boldsymbol m})}
       \kappa_{\boldsymbol m,\mathfrak a}
    =(c_0^2-c_0^{-r-s})\kappa_{\boldsymbol m},
 \quad r=m_1-1,\quad s=m_2-1,
\end{equation}
where
\[
 A_{\boldsymbol m}=i^{(r,s)}(\pi),\qquad
 B_{\boldsymbol m}=i^{(r,s)}(\bar\pi),\qquad
 A_{\boldsymbol m}B_{\boldsymbol m}=p^{r+s}.
\]
These algebraic numbers have complex absolute value
\(p^{(r+s)/2}>1\), because \(r,s\geq0\) and \(r+s>0\).
Both denominators and the numerator on the left are therefore
nonzero, including \(r=0\) and \(s=0\).
The right scalar is also nonzero.  The unit ambiguity in the
Artin action disappears since \(r-s\equiv0\pmod{w_K}\).
Thus \(\kappa_{\boldsymbol m}(h)=0\) for every
\(\boldsymbol m\in I\).
Corollary~\ref{cor:cm-radical-joint-kernel} now gives
\(h\bmod[H,H]\in N_{\mathrm{ell}}\).
This proves the exact equality
\(\ker\Phi_{c_0}=N_{\mathrm{ell}}\).

Finally \(G/\Gamma\) is finite of order prime to \(p\), so the
untwisted \(\mathbb Z_p[[G]]\) is finite free over \(\Lambda_2\).
The cyclotomic twist changes scalars by the continuous ring
automorphism \(\gamma\mapsto\chi_{\mathrm{cyc}}(\gamma)\gamma\)
of \(\Lambda_2\).  Its target remains torsion-free, and the
displayed injection proves the assertion.
\end{proof}

\begin{lemma}\label{lem:cm-exact-eigenlift-obstruction}
Let \(\gamma\neq1\) belong to \(\Gamma\), let
\(\lambda\in\mathbb Z_p^\times\), and let \(t\) be a lift of
\(\gamma\) to \(\operatorname{Gal}(\Omega/K)\).
If an element \(h\in H\) satisfies
\[
                         tht^{-1}=h^\lambda ,
\]
then every elliptic Soul\'e character vanishes on \(h\).
If \(\mathcal X_2^+\) is finite, then \(h\in[H,H]\).
\end{lemma}

\begin{proof}
Write \(a\) for the image of \(h\) in \(\overline A\).
The asserted group equation gives the exact module equation
\[
                              (\gamma-\lambda)a=0.
\]
The element \(\gamma-\lambda\) is nonzero in
\(\Lambda_2=\mathbb Z_p[[\mathbb Z_p^2]]\), an integral domain.
Lemma~\ref{lem:cm-mixed-measure-embedding} therefore gives \(a=0\).
The joint-kernel description implies the first assertion.
Under finite \(\mathcal X_2^+\),
Theorem~\ref{thm:cm-finite-plus-abelian} gives
\(N_{\mathrm{ell}}=0\), hence \(h\) has zero image in \(A\).
\end{proof}

This gives an obstruction to lifting homogeneous
Galois eigenclasses by group elements satisfying the corresponding
eigen-equations.  Rational eigenspace decompositions in a graded
quotient do not provide such equations in \(H\).

\begin{proposition}\label{prop:cm-no-normalized-cyclic}
Assume that \(\mathcal X_2^+\) is finite.  Put
\(G_1=\operatorname{Gal}(\Omega/L_1)\), so that
\[
             1\longrightarrow H\longrightarrow G_1
               \longrightarrow\Gamma\longrightarrow1 .
\]
If \(t\in G_1\) has nonidentity image in \(\Gamma\), then the
equation
\[
                 tht^{-1}=h^\lambda,\qquad
                    h\in H,\quad\lambda\in\mathbb Z_p^\times
\]
has only the solution \(h=1\).
Consequently, for every nontrivial closed procyclic subgroup
\(C\subset H\), the image of \(N_{G_1}(C)\) in \(\Gamma\) is
trivial.
\end{proposition}

\begin{proof}
The finite-plus theorem gives that \(H\) is free pro-\(p\) and
that \(A=H^{\mathrm{ab}}\) embeds equivariantly in
\[
 P=\mathbb Z_p[[G]](1).
\]
As a \(\Lambda_2\)-module, \(P\) is a finite direct sum of
copies of the regular module, with its cyclotomic twist.

We first justify taking completed tensor powers of this
injection.  Every compact torsion-free \(\mathbb Z_p\)-module
is topologically a product of copies of \(\mathbb Z_p\).
Indeed, its Pontryagin dual is a divisible \(p\)-primary group,
and hence a direct sum of copies of
\(\mathbb Q_p/\mathbb Z_p\).  For such a product \(M\), and any
compact \(\mathbb Z_p\)-module \(V\), the definition by finite
continuous quotients gives
\[
 \left(\prod_{i\in I}\mathbb Z_p\right)
       \widehat\otimes_{\mathbb Z_p}V
             \simeq\prod_{i\in I}V.
\]
This functor preserves injections.  Applying this assertion
one tensor factor at a time gives, for every \(d\geq1\), an
injection
\[
 A^{\widehat\otimes d}\lhook\joinrel\longrightarrow
 P^{\widehat\otimes d}
     \simeq\mathbb Z_p[[G^d]](d).
\]

The action of \(\Gamma\) on the last module is diagonal
translation with the scalar twist
\(\chi_{\mathrm{cyc}}^d\).  The continuous group isomorphism
\[
 \Gamma\times\Gamma^{d-1}\longrightarrow\Gamma^d,\qquad
 (g,h_2,\ldots,h_d)\longmapsto
                (g,gh_2,\ldots,gh_d)
\]
identifies diagonal translation with translation in the first
factor.  After choosing the finitely many tame cosets, its
completed group algebra is therefore a finite direct sum of
power-series modules
\[
             \Lambda_2[[T_1,\ldots,T_{2d-2}]].
\]
It is torsion-free over \(\Lambda_2\).  The scalar twist changes
the coefficient action by the ring automorphism
\(\gamma\mapsto\chi_{\mathrm{cyc}}(\gamma)^d\gamma\), which
preserves that property.  Thus every completed tensor power
of \(A\) is \(\Lambda_2\)-torsion-free.

Write \(\gamma_dH\) for the closed lower central series.
For a free pro-\(p\) group, the canonical map
\[
 \gamma_dH/\gamma_{d+1}H
        \longrightarrow A^{\widehat\otimes d},
 \qquad
 [h_1,\ldots,h_d]\longmapsto
          [\bar h_1,\ldots,\bar h_d]_{\mathrm{tensor}},
\]
identifies its source with the completed degree-\(d\) free
Lie submodule.  This assertion is integral: for a finite free
basis, the Hall commutators give the free
\(\mathbb Z_p\)-basis, and their Magnus expansions have the
corresponding leading Lie tensors.  For a general free
pro-\(p\) basis converging to one, projection onto finite
subsets and reduction modulo \(p^n\) give the completed
assertion.  These projections also prove injectivity, since a
nonzero coefficient is seen in a finite set of basis
coordinates and modulo some \(p^n\).  The construction is
functorial in automorphisms.  In particular, the conjugation
action on this graded piece is the action induced by that on
\(A\), and so it factors through \(\Gamma\).
It follows that every
\(\gamma_dH/\gamma_{d+1}H\) is
\(\Lambda_2\)-torsion-free.

Every pro-\(p\) group is separated for its closed lower central
series: any nonidentity element survives in a finite \(p\)-group,
which is nilpotent.  Therefore, if \(h\neq1\), it has a nonzero
image \(a\) in some
\(\gamma_dH/\gamma_{d+1}H\).  If \(t\) has image
\(\gamma\neq1\), the proposed equation gives
\[
                         (\gamma-\lambda)a=0 .
\]
The power series \(\gamma-\lambda\) is nonzero in
\(\Lambda_2\), contradicting the torsion-freeness just proved.
Thus \(h=1\).

Finally \(H\) is torsion-free, so a nontrivial closed
procyclic \(C\) is isomorphic to \(\mathbb Z_p\).
Conjugation by a normalizer element sends a topological
generator \(h\) to \(h^\lambda\) for a unit \(\lambda\).
The first assertion forces the image of that element in
\(\Gamma\) to be one.  This proves the precise normalizer
claim.
\end{proof}

The proposition gives a further constraint on a proposed
arithmetic use of the geometric recognition theorem.  A
nontrivial procyclic subgroup selected inside \(H\), including
one selected inside the arithmetic outer kernel, cannot
acquire an open normalizer image in
\(\operatorname{Gal}(L_\infty/L_1)\): its normalizer image
there is exactly the identity subgroup.  This conclusion
concerns the displayed arithmetic extension and follows from
its actual conjugation action on every lower-central quotient.

\paragraph{The remaining correction equation.}
Let \(\widetilde F^\bullet H\) be the weighted free filtration
attached to a proposed free set of arithmetic lifts, with its
filtered outer map to \(F^\bullet\rho(H)\).
A homogeneous derivation relation, such as the relation in
bidegree \((12,2)\) proved above, produces a word
\(r\in\widetilde F^{14}H\) satisfying
\[
 \rho(r)\in F^{15}\rho(H).
\]
To cancel that outer image without changing its initial
free class, one must solve the stronger, explicit equation
\[
 \rho(r)=\rho(s)\quad\hbox{for some }s\in\widetilde F^{15}H.
\]
Its solution gives \(rs^{-1}\in\ker\rho\).
The first inclusion alone does not supply \(s\):
it concerns the target filtration, whereas the displayed equation
asks for membership in the image of a specified source subgroup.
Rational surjectivity of successive graded maps gives no
uniform integral bound on the denominators in successive
corrections.  Clearing a growing denominator by successive
\(p\)-powers can erase the proposed initial class.
The argument in \cite[proof of Theorem 2.14, p.~27]{Ishii} uses graded
freeness precisely to identify these two filtrations.
The source-group freeness established here does not itself
solve the displayed correction equation.

\subsection{The exact mixed-chart image of the geometric measure}
\label{subsec:cm-measure-image}

Write \(w_K=|O_K^\times|\), put \(U=(\mathbb Z_p^\times)^2\subset V\), and let
\(j:U\hookrightarrow V\) be its inclusion.
The ray-class identification is
\(G=U/O_K^\times\).  For the auxiliary ideal
\(\mathfrak a=(c)\), let \(\Phi_c\) denote the map of Lemma~\ref{lem:cm-mixed-measure-embedding} composed with \(A\to\overline A\), and let \(\mu_c\) be its normalized
\(O_K^\times\)-invariant lift of the mixed-unit
measure \(\Phi_c\) to \(U\): its mass on each finite
unit orbit is the quotient mass divided by \(w_K\).
This is integral since \(p\nmid w_K\).
Let \(P,Q\) act on \(V\) by multiplication by
\(\pi,\bar\pi\), and put
\[
 \mathcal L=(1-P_*)(1-Q_*),\qquad
 D_c=c^2-[c^{-1}]_*.
\]
Here all starred maps denote pushforward of bounded
measures, and \(PQ=p\,\operatorname{id}_V\).

\begin{lemma}\label{lem:cm-measure-euler-image}
The actual geometric measure satisfies
\[
 w_Kj_*\mu_c=D_c\mathcal L\alpha,\qquad
 w_Kj_*\mu_c=D_cj_*(\alpha|_U),\qquad
 \mathcal L\alpha=j_*(\alpha|_U).
\]
Its restriction \(\alpha|_U\) has kernel
\(N_{\mathrm{ell}}\), and uniquely determines
\(\alpha\) with its zero-total-mass normalization.
Under finite \(\mathcal X_2^+\), this restriction
embeds \(A\) as a full-rank closed submodule of
\(\mathbb Z_p[[U]]^{O_K^\times}(1)\).
After identifying this target with
\(\mathbb Z_p[[G]](1)\) and absorbing the twist,
the image is an ideal containing a nonzerodivisor.
\end{lemma}

\begin{proof}
The \((r,s)\)-moment of \(\mu_c\) is
\(\kappa_{(r+1,s+1),\mathfrak a}\).  Rearranging
\cite[Proposition 4.3, pp.~17--18]{IshiiSoule}
and using \cite[Theorem 2.6, p.~8]{Ishii}
gives exactly the first identity on every admissible
positive-degree monomial.  Inadmissible moments
vanish by \(O_K^\times\)-invariance.  The constant
moment of \(\mu_c\) is zero: it is the Kummer value
of \(N_{L_n/K}(u_n)\), a global unit of \(K\), hence
a root of unity with all \(p\)-power roots in
\(L_\infty\).  The constant moment of \(\alpha\)
is zero by its source formula, and the odd moments
vanish by the action of \(-1\).
Polynomial moments determine bounded measures,
so the first identity holds integrally.

Neither \(P(V)\) nor \(Q(V)\) meets \(U\).
Restriction to \(U\), which commutes with \(D_c\),
therefore gives the second identity.
The \((r,s)\)-moment multiplier of \(D_c\) is
\(c^2-c^{-r-s}\ne0\), including degree zero.
Thus \(D_c\) is injective, and cancellation gives
the third identity.

The kernel of \(\mathcal L\) on all bounded
measures is precisely \(\mathbb Z_p\delta_0\).
Indeed its positive-degree moment multipliers
are
\[
 (1-i^{(r,s)}(\pi))(1-i^{(r,s)}(\bar\pi)),
\]
and both algebraic numbers occurring here have
complex absolute value \(p^{(r+s)/2}>1\).
A measure killed by \(\mathcal L\) consequently
has all positive-degree moments zero and is its
total mass times \(\delta_0\).
This proves uniqueness with zero total mass.
The second identity and injectivity of \(D_c\)
identify \(\ker(\alpha|_U)\) with
\(\ker\Phi_c=N_{\mathrm{ell}}\).

Under finite \(\mathcal X_2^+\), the latter kernel
vanishes.  The exact rank formula
\cite[Theorem 4.4, author version p.~27]{VenjakobIwasawa},
together with weak Leopoldt, gives
\[
 \operatorname{rank}_{\Lambda_2}A
     =r_2(L_1)=[L_1:K]=\frac{(p-1)^2}{w_K}.
\]
The maximal pro-\(p\) group comparison used here
is the \((p,S)\)-closed comparison at the start of
\cite[Section 4, author version p.~26]{VenjakobIwasawa}.
The unit-chart target has this same rank.
Equivariance makes the image an ideal after the
cyclotomic twist is absorbed.  The tame
decomposition of \(\mathbb Z_p[[G]]\) is a product
of copies of \(\Lambda_2\), and full rank makes
each ideal component nonzero.  A choice of a
nonzero element in every component gives the
claimed nonzerodivisor.
\end{proof}

These identities retain the radial and axis constraints
on the actual image.  In particular, for \(i,j\geq0\),
\[
 \alpha|_{P^iQ^jU}=(P^iQ^j)_*(\alpha|_U).
\]
This follows by restricting the Euler identity
successively to the disjoint interior shells.
No infinite geometric-series expansion at the axes
is required.  In the class-two question, the remaining
map is still
\[
 \mathcal J\circ(\alpha\widehat\otimes\alpha)
              \bigm|_{\widehat{\bigwedge}^{\,2}A}.
\]
The separate-factor operators \(\mathcal L\otimes1\)
and \(1\otimes\mathcal L\) do not commute with its
addition shears.  The displayed ideal description
therefore retains, rather than removes, this exact
tensor-image constraint.

\subsection{The finite elliptic Euler complexes and the actual kernel}
\label{subsec:cm-elliptic-euler-kernel}

There is a finite-level description of the precise arithmetic
image entering \(J_A\).  Retain the finite-\(\mathcal X_2^+\)
hypothesis and write
\[
 V_n=E[p^n],\qquad
 \alpha_n:A\longrightarrow\mathbb Z_p[V_n](1),\qquad
 A_n=\operatorname{im}\alpha_n.
\]
The transition on \(V_n\) is multiplication by \(p\).
Each \(A_n\) is its actual compact, finite free
\(\mathbb Z_p\)-lattice.

Use Goncharov's elliptic units \(\theta_E(x)\), and let
\(C_1(n)\) be their \(\mathbb Q_p\)-span for
\(0\ne x\in V_n\), inside \(L_n^\times\otimes\mathbb Q_p\).
Here \(B_2(F)\) denotes Goncharov's rational Bloch group.
It is generated by symbols \(\{z\}\), \(z\in\mathbb P^1(F)\),
with differential \(\delta\{z\}=(1-z)\wedge z\), set equal to
zero at \(0,1,\infty\). Its relations identify the two
specializations of \(\sum_i\{f_i\}\) along a connected curve
whenever \(\sum_i(1-f_i)\wedge f_i=0\) in its function field;
see \cite[Section 2.1, p.~10]{GoncharovEuler}.
For number fields this agrees rationally with the five-term
presentation through the canonical map stated there.
Let \(C_2(n)\) be the \(\mathbb Q_p\)-span of his
double elliptic units \(\theta_E(a,b,c)\), with
\(a,b,c\in V_n\) and \(a+b+c=0\), in
\(B_2(L_n)\otimes\mathbb Q_p\).
The ray field and these theta-value spans are described in
\cite[Definition 5.2, Proposition 5.3 and Definition 5.4,
p.~27]{GoncharovEuler}.
Its differential is
\begin{equation}\label{eq:cm-elliptic-triangle-delta}
 \delta\theta_E(a,b,c)
   =\theta_E(a)\wedge\theta_E(b)
    +\theta_E(b)\wedge\theta_E(c)
    +\theta_E(c)\wedge\theta_E(a),
\end{equation}
by \cite[Theorem 2.10(a), equation (38), p.~15]{GoncharovEuler}.
The differential of a triple with a zero entry vanishes;
no value of \(\theta_E(0)\) is needed in its image.
Put
\[
 \mathcal R_n^{\rm ell}
   =\operatorname{coker}\bigl(C_2(n)
                          \xrightarrow{\delta}\bigwedge^2C_1(n)\bigr),
 \qquad
 \mathcal D_n=\ker\bigl(J_n:\bigwedge^2A_n
                   \longrightarrow\mathbb Z_p[V_n\times V_n](2)\bigr),
\]
where \(J_n\) is the same pair of finite additive shears
as in Theorem~\ref{thm:cm-measure-class-two}.

\begin{theorem}\label{thm:cm-elliptic-euler-kernel}
There are natural equivariant identifications
\[
 A_n[1/p]=\operatorname{Hom}_{\mathbb Q_p}
                         (C_1(n),\mathbb Q_p(1)),
 \qquad
 \mathcal D_n[1/p]=\operatorname{Hom}_{\mathbb Q_p}
                         (\mathcal R_n^{\rm ell},\mathbb Q_p(2)).
\]
The transition \(\mathcal D_{n+1}\to\mathcal D_n\)
is induced by the actual pushforward
\(A_{n+1}\to A_n\).  Rationally, it is dual to the
natural map \(\mathcal R_n^{\rm ell}\to
\mathcal R_{n+1}^{\rm ell}\) induced by field extension
and the elliptic-unit distributions.  Finally,
\begin{equation}\label{eq:cm-J-kernel-elliptic-limit}
                       \ker J_A=\varprojlim_n\mathcal D_n.
\end{equation}
Here the inverse limit retains the specified integral
lattices \(\mathcal D_n\).
\end{theorem}

\begin{proof}
We first check the normalization linking the two primary
sources.  Scale the period lattice to
\(L=\mathbb Z\tau+\mathbb Z\), and write
\(z=a\tau+b\), \(q=e^{2\pi i\tau}\),
\(\xi=e^{2\pi iz}\).  Put
\(\eta_1=\eta(\tau)\), \(\eta_2=\eta(1)\).
The sigma product and the Legendre relation give
\[
 \sigma(z,L)=\frac{e^{\eta_2z^2/2}}{2\pi i}
       \xi^{-1/2}(\xi-1)
       \prod_{m\ge1}
       \frac{(1-q^m\xi)(1-q^m\xi^{-1})}{(1-q^m)^2},
 \qquad \eta_2\tau-\eta_1=2\pi i.
\]
The product identity follows from the zeros and the two
quasiperiods of \(\sigma\), with \(\sigma'(0)=1\)
fixing its constant.  Substitution in Nakamura's
\(\Theta(z,L)=\Delta(L)e^{-6\eta(z)z}\sigma(z,L)^{12}\),
using \(\eta(z)=a\eta_1+b\eta_2\), yields
\[
 \Theta(z,L)=q^{6a^2-6a+1}e^{12\pi ib(a-1)}
    (1-\xi)^{12}
    \prod_{m\ge1}(1-q^m\xi)^{12}(1-q^m\xi^{-1})^{12}.
\]
This is precisely the twelfth power of
\cite[equation (19), p.~7]{GoncharovEuler}.
Homogeneity removes the scaling of the lattice.
The root-of-unity ambiguity from torsion representatives
has zero Kummer character on \(H\).  Thus
\(\kappa_{\Theta(x)}=12\kappa_{\theta_E(x)}\);
the factor \(12\) is invertible because \(p\ge5\).

Nakamura's formulas are stronger than a moment identity.
By \cite[equation (3.5.1), p.~209, and
equations (4.11.3)--(4.11.5), p.~220]{Nakamura1995},
the actual finite coefficients satisfy
\begin{equation}\label{eq:cm-alpha-theta-finite}
 \alpha_n(h)(x)=12\kappa_{\theta_E(x)}(h)\quad(x\ne0),
 \qquad
 \alpha_n(h)(0)=-n\kappa_{p^{12}}(h).
\end{equation}
The latter formula also uses
\cite[Proposition (3.8), p.~210]{Nakamura1995}.
Ishii's special-lift measure is this measure, with basis
compatibility recorded in
\cite[Remark 2.6.1(2), p.~8]{Ishii}.
CM-unit invariance and the ray-field description in
\cite[p.~27]{GoncharovEuler} place these rational theta
classes over \(L_n\).

We verify that their restriction to \(H\) loses no
nonzero rational class.  Set
\(G_n=\operatorname{Gal}(\Omega/L_n)\) and
\(\Gamma_n=\operatorname{Gal}(L_\infty/L_n)\).
Finite-field Kummer theory injects the \(p\)-unit
space of \(L_n\) into \(H^1(G_n,\mathbb Q_p(1))\):
the relevant Kummer extensions are pro-\(p\) and
unramified outside \(p\), and hence lie in \(\Omega\).
The restriction kernel to \(H\) is the image of
\(H^1(\Gamma_n,\mathbb Q_p(1))\), which is zero.
Indeed, choose \(\gamma_0\) with
\(\chi_{\rm cyc}(\gamma_0)\ne1\).  Since \(\Gamma_n\)
is abelian, every cocycle \(f\) satisfies
\[
 (\chi_{\rm cyc}(\gamma_0)-1)f(\gamma)
       =(\chi_{\rm cyc}(\gamma)-1)f(\gamma_0),
\]
so it is a coboundary over \(\mathbb Q_p\).
This proves the required Kummer injectivity, including
the possible infinitely \(p\)-divisible kernel in the
base tower.

The nonzero coordinate functionals in
\eqref{eq:cm-alpha-theta-finite} span precisely
\(C_1(n)\); the origin coordinate is their negative
sum, since \(\alpha\) has total mass zero.
A proper rational image of \(A_n[1/p]\) in the dual
of \(C_1(n)\) would have a nonzero annihilator there,
contradicting the Kummer injectivity just proved.
This proves the first identification, while retaining
\(A_n\) as its actual integral lattice.

For \(\eta\in\bigwedge^2A_n[1/p]\), let
\(B_\eta\) denote the alternating pairing on \(C_1(n)\)
under that identification.  The inverse-image formula
for the shears is
\[
 (J_n\eta)(x,y)=\eta(x,y-x)+\eta(x-y,y)-\eta(x,y).
\]
When \(x,y,x-y\ne0\), evenness
\(\theta_E(-u)=\theta_E(u)\) and
\eqref{eq:cm-alpha-theta-finite} give
\begin{equation}\label{eq:cm-J-elliptic-triangle}
 (J_n\eta)(x,y)
       =144 B_\eta\bigl(\delta\theta_E(x,y-x,-y)\bigr).
\end{equation}
When \(x=0\), \(y=0\), or \(x=y\), the coefficient
vanishes by alternation and evenness, independently
of the origin coefficient.  Every nondegenerate
zero-sum triple occurs in
\eqref{eq:cm-J-elliptic-triangle}.
Thus the rational kernel is exactly the dual of
\(\mathcal R_n^{\rm ell}\), with the second Tate
twist.  Its integral kernel \(\mathcal D_n\) is
closed and \(p\)-saturated because its target is
finite free over \(\mathbb Z_p\).

Here are the transition maps, including the origin.
For \(x\ne0\), the isogeny distribution gives, in
additive notation for units tensored with \(\mathbb Q_p\),
\[
 \theta_E(x)=\sum_{[p]y=x}\theta_E(y).
\]
This is the ordinary inclusion of the theta class at
the same torsion point into the next field.  Hence
pushforward \(A_{n+1}\to A_n\) is dual to
\(C_1(n)\hookrightarrow C_1(n+1)\).
At zero the exact calculation is
\[
 \sum_{[p]y=0}\alpha_{n+1}(h)(y)
   =-(n+1)\kappa_{p^{12}}(h)
         +\sum_{0\ne y\in E[p]}\kappa_{\Theta(y)}(h)
   =-n\kappa_{p^{12}}(h),
\]
using \cite[Proposition (2.2)(3), p.~204]{Nakamura1995}.
For double units, independence of the torsion exponent
and the sum over preimage triples are
\cite[Lemma 2.12, p.~17, and equation (41), p.~16]{GoncharovEuler}.
Their differential is the exterior square of the
unit inclusion: two chosen preimages determine the
third by the zero-sum condition.  This proves the
commuting diagram of complexes and the asserted
natural map on cokernels.

Finally, compactness and injectivity of \(\alpha\)
identify \(A\) with \(\varprojlim A_n\).
The coefficient quotients modulo \(p^r\), for all
\(n,r\), form a cofinal system of neighborhoods.
Completed tensor and exterior squares therefore give
\(\widehat{\bigwedge}^2A=\varprojlim\bigwedge^2A_n\).
The shear maps commute with these projections,
so taking their kernels proves
\eqref{eq:cm-J-kernel-elliptic-limit}.
\end{proof}

The transition just used has an explicit comparison
with arithmetic norms.  On a torsion point whose two
Tate coordinates are units,
\(\operatorname{Gal}(L_{n+1}/L_n)\) acts by
\((1+p^nu,1+p^nv)\), \(u,v\in\mathbb F_p\), and
its orbit is the entire \(p^2\)-element \([p]\)-fiber.
The distribution product is then the field norm.
If one coordinate is divisible by \(p\) and the other
is a unit, the arithmetic orbit has \(p\) elements,
each repeated \(p\) times in the norm, whereas the
isogeny fiber has \(p^2\) elements.  The all-chart
transition in the theorem is the isogeny distribution,
with the displayed correction at zero.

For the Gaussian and Eisenstein fields,
\cite[Theorem 5.16, p.~34]{GoncharovEuler} supplies a
surjective morphism from the Bianchi modular complex
at the ideal \((p^n)\) to the complex with target
\(\bigwedge^2\widehat C_1(n)\), where
\(\widehat C_1(n)=C_1(n)\oplus\mathbb Q_p\theta_E(0)\).
Compose its degree-two map with the exterior square
of the projection \(\widehat C_1(n)\to C_1(n)\)
that sends the formal \(\theta_E(0)\) to zero.
The differentials of degenerate triangles vanish,
so this is a surjective morphism to the exact Euler
complex used above.  One can equally send the formal
origin to its actual normalized class \(-n[p]\).
At a prime ideal \(\mathfrak P\),
\cite[Theorem 1.2, p.~9]{GoncharovEuler} identifies
\[
 \operatorname{coker}\bigl(C_2(\mathfrak P)
       \longrightarrow\bigwedge^2C_1^{\rm un}(\mathfrak P)\bigr)
       =H^2_{\rm cusp}(\Gamma_1(\mathfrak P),\mathbb Q),
\]
where \(C_1^{\rm un}\) is the global-unit space.
Since the differential lands in that smaller exterior
square, this gives a canonical subspace of the full
unit cokernel after extension of scalars.  The
global-unit qualification is part of this precise map.

The remaining persistence issue has a concrete form.
Injectivity of the natural map on cokernels requires
\[
 \delta C_2(n+1)\cap\bigwedge^2C_1(n)=\delta C_2(n).
\]
The inclusions of unit spaces alone do not establish
this equality.  Moreover,
\eqref{eq:cm-J-kernel-elliptic-limit} retains the
integral lattices, so rational finite-level classes
must have compatible bounded duals to contribute to
the actual kernel.  The theorem supplies this exact
finite-level and inverse-limit comparison without
inferring such persistence from a dimension count.

\subsection{The centralizer pullback and its kernel}
\label{subsec:cm-actual-centralizer}

Let \(F\) be the geometric free pro-\(p\) group of the
once-punctured CM elliptic curve. The prime-to-\(p\) degree
of \(K(E[p])/L_1\) gives
\(K(E[p])\cap\Omega=L_1\), and restriction identifies
\(\operatorname{Gal}(K(E[p])\Omega/K(E[p]))\) with
\(G_1=\operatorname{Gal}(\Omega/L_1)\).
Ishii's kernel-field inclusion therefore defines the map
\[
 \rho:G_1\longrightarrow\operatorname{Out}(F)
\]
used here.
Put \(O=\rho(G_1)\), \(K=\ker\rho\), and let
\(\widetilde O\) be the inverse image of \(O\) in
\(\operatorname{Aut}(F)\).  The following pullback records the
arithmetic kernel without making a choice of automorphism lifts:
\[
 \mathcal P=\widetilde O\times_O G_1
    =\{(\alpha,g):[\alpha]=\rho(g)\}.
\]
Its geometric subgroup is
\(F\longrightarrow\mathcal P,\ f\longmapsto(\operatorname{inn}(f),1)\).
This is injective because \(F\) has trivial center.

\begin{proposition}\label{prop:cm-actual-centralizer-pullback}
There are exact identifications
\[
 C_{\mathcal P}(F)=\{(1,k):k\in K\}\simeq K,
 \qquad
 \mathcal P|_K=F\times K.
\]
For every closed subgroup \(A\subset F\),
\[
 N_{\mathcal P}(A)=N_{\widetilde O}(A)\times_O G_1.
\]
In particular, if \(V_A\subset O\) is the image of
\(N_{\widetilde O}(A)\), then the arithmetic normalizer image
is precisely \(\rho^{-1}(V_A)\).  It is open in \(G_1\) exactly
when \(V_A\) is open in \(O\).
\end{proposition}

\begin{proof}
Conjugation by \((\alpha,g)\) on the geometric subgroup is
the automorphism \(\alpha\).  It centralizes all of \(F\)
exactly when \(\alpha=1\), which forces \(g\in K\).
Above \(K\), every \(\alpha\) is inner.  The unique element
of \(F\) inducing it gives the displayed direct product,
with multiplication inherited from the pullback.

The same conjugation formula shows that an element normalizes
\(A\) exactly when \(\alpha(A)=A\).  This proves the
normalizer identity and its projection formula.  Finally
the continuous surjection \(G_1\to O\) has the quotient
topology, so a subgroup of \(O\) is open exactly when
its inverse image is open.
\end{proof}

Thus the geometric open-normalizer test factors through the
outer image \(O\).  To prove arithmetic faithfulness
from this pullback one must prove the further assertion
\(C_{\mathcal P}(F)=1\), or equivalently \(K=1\).
The proposition gives the exact map at that step.

There is also an arithmetic restriction on possible mixed
procyclic subgroups of the pullback.  Assume
\(\mathcal X_2^+\) is finite.  If \(A\subset F\times K\)
is closed procyclic and its projection to \(K\) is nontrivial,
then its normalizer in \(\mathcal P\) has trivial image
in \(\Gamma\).  Indeed the projection of any normalizer
element normalizes the nontrivial procyclic image of \(A\)
inside \(K\subset H\), so
Proposition~\ref{prop:cm-no-normalized-cyclic} applies.
Consequently a procyclic subgroup of \(F\times K\)
with open \(\Gamma\)-image normalizer must lie in \(F\).
This conclusion uses the centralizer pullback and
the arithmetic conjugation modules; it leaves the possible
nontrivial, infinitely generated kernel \(K\) as the
remaining centralizer question.

For completeness, the asserted infinite generation also
follows from those same modules.  If a nontrivial finitely
generated closed subgroup \(J\subset H\) were normalized
by a lift of \(\gamma\neq1\), choose the first nonzero
image of \(J\) in \(\gamma_dH/\gamma_{d+1}H\).
It is a finitely generated torsion-free
\(\mathbb Z_p\)-module preserved by \(\gamma\).
Its characteristic polynomial would be a nonzero
element of \(\Lambda_2\) annihilating that image,
contrary to the proved torsion-freeness of the
lower-central module.  Apply this to the normal subgroup
\(K\) to obtain the stated conclusion whenever \(K\neq1\).


\subsection{The examples over \(\mathbb Q(i)\) and
\(\mathbb Q(\sqrt{-3})\)}

We verify both arithmetic hypotheses of Ishii's theorem by
finite calculations, leaving its freeness hypothesis explicit.
The cases are those recorded in
\cite[Example 1.2, p.~3]{Ishii}.

\begin{proposition}\label{prop:cm-small-primes}
For \((K,p)=(\mathbb Q(i),5)\) and
\((\mathbb Q(\sqrt{-3}),7)\), the ray class field \(K(p)\) has
class number prime to \(p\), and \(K(p^2)\) has exactly two primes
above \(p\). For the reductions of
\(E:y^2=x^3-x\) at \(5\) and \(E:y^2=x^3+1\) at \(7\),
the Frobenius polynomials are respectively
\[
 U^2+2U+5,\qquad U^2+4U+7 .
\]
\end{proposition}
\begin{proof}
For a class-number-one imaginary quadratic field, the ray class
sequence gives
\[
 \operatorname{Gal}(K(p^n)/K)
  \simeq(\mathcal O_K/p^n)^\times/\mathcal O_K^\times .
\]
The units inject modulo \(p\), and the Chinese remainder map at
the two split primes gives degree
\((p^{n-1}(p-1))^2/w_K\).
At \(n=1\) this is \(4\) and \(6\), respectively.
The cyclotomic extension \(K(\mu_p)/K\) has degree \(p-1\) and
conductor dividing \(p\); hence
\[
 \mathbb Q(i)(5)=\mathbb Q(\zeta_{20}),\qquad
 \mathbb Q(\sqrt{-3})(7)=\mathbb Q(\zeta_{21}) .
\]
Equivalently, at level one the norm
\((u,v)\mapsto uv\) identifies
\((\mathbb F_p^\times)^2/\mathcal O_K^\times\)
with \(\mathbb F_p^\times\), since \(w_K=p-1\) in these two cases.

For \(K=\mathbb Q(i)\), take \(\pi=2+i\), \(\bar\pi=2-i\).
The isomorphism \(\mathcal O_K/\bar\pi^2\simeq\mathbb Z/25\)
sends \(i\) to \(7\). It sends \(\pi\) to \(9\), and
\(9^4\equiv11\not\equiv1\pmod {25}\).
The quotient
\((\mathcal O_K/\bar\pi^2)^\times/\mathcal O_K^\times\)
has order \(5\); the class of \(\pi\) generates it.
For \(K=\mathbb Q(\sqrt{-3})\), put
\(\omega=(-1+\sqrt{-3})/2\), \(\pi=3+\omega\), and
\(\bar\pi=2-\omega\). The quotient by \(\bar\pi^2\)
sends \(\omega\) to \(30\pmod {49}\), so \(\pi\) maps to \(33\).
Here \(33^6\equiv8\not\equiv1\pmod {49}\), and the corresponding
unit quotient has order \(7\).
In each case inertia at \(\mathfrak p\) supplies the first local
unit factor; the Frobenius at \(\mathfrak p\) supplies the generator
just calculated of the remaining quotient. The decomposition group
is therefore the full ray class group. Each of the two primes of
\(K\) above \(p\) has one extension to \(K(p^2)\).
This is the explicit reciprocity criterion in
\cite[Remark 2.14.1, p.~11]{Ishii}.

We next verify the prime-to-\(p\) class numbers. Minkowski's bound
provides representatives for every ideal class of norm at most
\[
 B_L=(4/\pi)^{r_2}\frac{[L:\mathbb Q]!}
 {[L:\mathbb Q]^{[L:\mathbb Q]}}\sqrt{|d_L|}
\]
\cite[Theorem 4.3, p.~70]{MilneANT}.
For \(L=\mathbb Q(\zeta_{20})\), its degree is \(8\),
\(|d_L|=2^8 5^6\), and \(\sqrt{|d_L|}=2000\).
The elementary estimate \(\pi>3\) gives
\(B_L<4375/288<16\).
The residue degree at \(2\) is \(4\), so the prime norms above
\(2\) equal \(16\). At \(3,7,11,13\) the residue degrees are
\(4,4,2,4\); their prime norms exceed the bound.
Only the primes above \(5\), of norm \(5\), can generate the
class group. The two primes above \(5\) have ramification index
\(4\) over \(\mathbb Q(i)\), and their fourth powers are the
principal ideals \((2+i)\) and \((2-i)\). The class group is
therefore generated by elements killed by \(4\).
Its order is prime to \(5\).

For \(L=\mathbb Q(\zeta_{21})\), its degree is \(12\),
\(|d_L|=3^6 7^{10}\), and \(\sqrt{|d_L|}=453789\).
Using \(\pi>31/10\) gives
\[
 B_L<
 \frac{8088368750000}{71887798161}<113.
\]
At \(2\) the residue degree is \(6\), giving norm \(64\).
At \(3\) the residue degree is \(6\), giving norm \(729\).
The primes above \(7\) have norm \(7\). Among unramified primes
with residue degree one below \(113\), the congruence
\(q\equiv1\pmod {21}\) leaves only \(q=43\).
For residue degree at least two, \(q^2<113\) restricts \(q\)
to \(2,3,5,7\), and \(5\) has residue degree \(6\).
Thus the primes of possible norm below the bound lie over
\(2,7,43\).

Let \(\xi=\zeta_7\) and
\(\alpha=\xi+\xi^2+\xi^4\). The identity
\(\alpha^2+\alpha+2=0\) shows that
\(|N_{L/\mathbb Q}(\alpha)|=64\).
Thus \((\alpha)\) is one of the primes above \(2\), and conjugation
shows that all of them are principal.
Let \(\zeta=\zeta_{21}\) and \(\beta=1+\zeta-\zeta^3\).
Its norm is \(43\). One direct computation is
\[
 \Phi_{21}(u)\equiv3u^2+4u+4\pmod {u^3-u-1},
 \qquad
 \operatorname{Res}(u^3-u-1,3u^2+4u+4)=43 .
\]
Consequently the primes above \(43\) are also principal.
The primes above \(7\) have sixth powers \((3+\omega)\) and
\((2-\omega)\). The class group is generated by elements killed
by \(6\); its order is prime to \(7\).

Finally, for \(y^2=x^3-x\) over \(\mathbb F_5\), the numbers of
solutions in \(y\) for \(x=0,1,2,3,4\) are
\(1,1,2,2,1\). Adding the point at infinity gives \(8\) points
and trace \(5+1-8=-2\).
For \(y^2=x^3+1\) over \(\mathbb F_7\), the corresponding numbers
are \(2,2,2,1,2,1,1\). There are \(12\) points and trace
\(7+1-12=-4\). These prove the two displayed polynomials for
the Frobenius endomorphisms of the reductions.
\end{proof}

\subsection{Locality in the two examples, with the unit calculation made explicit}

\begin{proposition}\label{prop:cm-locality}
For the two pairs of Proposition~\ref{prop:cm-small-primes},
put \(L=K(\mu_p)=K(p)\), and fix a prime \(\mathfrak p\) above \(p\).
The decomposition homomorphism is an isomorphism
\[
 G_{\mathbb Q_p(\mu_p)}^{(p)}
    \xrightarrow{\ \sim\ }\operatorname{Gal}(\Omega/L).
\]
\end{proposition}
\begin{proof}
Put \(\Delta=\operatorname{Gal}(L/K)\), and identify its
cyclotomic character with \(\omega\). Let
\(V=\mathcal O_L[1/p]^\times/p\). Kummer theory and
\(p\nmid h(L)\) identify \(V\) with \(H^1(G,\mu_p)\),
where \(G=\operatorname{Gal}(\Omega/L)\).
The two localization maps
\[
 V\longrightarrow
 L_{\mathfrak p}^{\times}/p\ \oplus\
 L_{\bar{\mathfrak p}}^{\times}/p
\]
are jointly injective. Indeed a unit in their joint
kernel would give a degree-\(p\) Kummer extension that
is split at the two primes above \(p\) and unramified
elsewhere, contradicting \(p\nmid h(L)\).

Complex conjugation exchanges the two primes.
Consequently localization at \(\mathfrak p\) alone is
injective on the image of
\(\mathcal O_{L^+}[1/p]^\times/p\).
The extension map from this real unit group is injective,
since norm followed by extension is multiplication by two.
Its dimension is \(p-1\), and its \(\Delta\)-representation
is the regular representation. In contrast \(V\) has
dimension \(p+1\) and representation
\[
 V\simeq\mathbb F_p[\Delta]\oplus\mathbb F_p
                      \oplus\mathbb F_p(1).
\]
The first two terms describe the free \(S\)-units, and the
last the \(p\)-power roots of unity. Hence the kernel of
localization at \(\mathfrak p\), being \(\Delta\)-stable,
can occur only in the characters \(1\) and \(\omega\).

The trivial-character component is generated by the
two units \(\pi,\bar\pi\) from \(K\). At \(\mathfrak p\),
\(\pi\) has valuation one and \(\bar\pi\) has nonzero
principal-unit class in \(\mathbb Q_p^\times/
(\mathbb Q_p^\times)^p\). The latter assertion is exactly
the conjugate of the congruence modulo \(p^2\) proved in
Proposition~\ref{prop:cm-small-primes}. The two local
classes are independent. They remain so after adjoining
\(\mu_p\), since that local extension has degree \(p-1\)
and its norm proves injectivity on classes modulo \(p\).
Thus the trivial-character part of the kernel is zero.

The \(\omega\)-component of \(V\) corresponds, after
Kummer untwisting, to \(H^1(G,\mathbb F_p)^\Delta\).
This has dimension two and is the inflation of
\(H^1(\Gamma,\mathbb F_p)\), where
\(\Gamma=\operatorname{Gal}(\bigcup_nK(p^n)/L)\).
The decomposition group in \(\Gamma\) is all of \(\Gamma\):
its image modulo \(p\) is already full by the calculation
at modulus \(p^2\), and a closed subgroup of \(\mathbb Z_p^2\)
surjecting modulo \(p\) is the whole group. Therefore
restriction to the local Galois group is injective on
this last character component as well.
We have proved that the whole localization map on \(V\)
is injective.

The local-to-global homomorphism \(Q\to G\), with
\(Q=G_{\mathbb Q_p(\mu_p)}^{(p)}\), is consequently
surjective: a proper closed image in a pro-\(p\) group
would be annihilated by a nonzero homomorphism to
\(\mathbb F_p\). Both \(H^1(G,\mathbb F_p)\) and
\(H^1(Q,\mathbb F_p)\) have dimension \(p+1\), so the
induced injection is an isomorphism.

Restriction \(H^2(G,\mu_p)\to H^2(Q,\mu_p)\) is
injective. The global group injects into the
\(p\)-torsion of the \(S\)-integer Brauer group, identified
by local invariants with
\(\{(a,b)\in\mathbb F_p^2:a+b=0\}\); localization at
\(\mathfrak p\) is projection to the first coordinate
and is injective on this group. These maps commute
with restriction in Galois cohomology.

For \(N=\ker(Q\to G)\), the five-term sequence now
gives \(H^1(N,\mathbb F_p)^G=0\).
A nontrivial pro-\(p\) group \(N\) has a nonzero
homomorphism to \(\mathbb F_p\). Its finite orbit
under the continuous pro-\(p\) action spans a nonzero
finite-dimensional representation of a finite
\(p\)-group, which has a nonzero fixed vector.
Thus \(H^1(N,\mathbb F_p)^G\ne0\) whenever \(N\ne1\).
The contradiction proves \(N=1\).
\end{proof}

The real \(S\)-unit dimension and the extra trivial-character
component in this proof repair the corresponding dimension
step in the proof of \cite[Theorem A.1, pp.~28--29]{Ishii}.
For these examples the numerical dimensions are \(4,6\) at
\(p=5\), and \(6,8\) at \(p=7\), for the real and complex
groups respectively.

For completeness, this also repairs the locality input to
nonvanishing of all the admissible elliptic Soul\'e characters
in these examples. Here \(w_K=p-1\), so an admissible pair
satisfies \(m_1\equiv m_2\pmod{p-1}\). The local character is
\(\chi_{\mathrm{cyc}}^{m_1}\eta^{m_2-m_1}\), where \(\eta\)
is the unramified CM unit-root character. Its Frobenius
value is an algebraic number of complex absolute value
\(\sqrt p\), and hence is not a root of unity.
Local Tate duality gives
\[
 H^2(\mathbb Q_p(\mu_p),\mathbb Q_p(\boldsymbol m))=0
 \qquad(\boldsymbol m\in I_K):
\]
the dual representation has an invariant only when inertia
forces \(m_1=1\) and the unramified character then forces
\(m_2=1\), the excluded pair. Locality and prime-to-\(p\)
descent give global \(H^2\) finiteness. The nonvanishing
assertion follows from \cite[Theorem 2.8(2), p.~9]{Ishii}.

\subsection{An unconditional abelian quotient for the two CM curves}

Let \(\Omega^*\subset\Omega\) be the pro-\(p\) field constructed
in the proof of Ishii's Lemma 2.13, p.~10. Its proved field
identity is
\[
 \overline K^{\ker\rho_{X,p}}=K(E[p])\Omega^* .
\]
The finite extension \(K(E[p])/K(p)\) has degree prime to \(p\),
and its factor in the outer image is central. The action on
the CM Tate module gives \(L_\infty=\bigcup_nK(p^n)\subset\Omega^*\).
Thus the kernel of \(G=\operatorname{Gal}(\Omega/K(p))
\to\operatorname{Gal}(\Omega^*/K(p))\) is contained in
\(H=\operatorname{Gal}(\Omega/L_\infty)\).

\begin{proposition}\label{prop:cm-abelian-detection}
For \(E:y^2=x^3-x\) over \(K=\mathbb Q(i)\), \(p=5\), and
for \(E:y^2=x^3+1\) over \(K=\mathbb Q(\sqrt{-3})\), \(p=7\),
the joint elliptic Soul\'e map is injective:
\[
 A=H^{\mathrm{ab}}
   \longrightarrow\prod_{\boldsymbol m\in I_K}
                 \mathbb Z_p(\boldsymbol m),
 \qquad a\longmapsto(\kappa_{\boldsymbol m}(a))_{\boldsymbol m}.
\]
In particular,
\[
 \operatorname{Gal}(\Omega/\Omega^*)\subset
       \overline{[H,H]},
 \qquad
 \Omega^{\overline{[H,H]}}\subset\Omega^* .
\]
\end{proposition}
\begin{proof}
The arithmetic hypotheses used in this proof have been
established in Proposition~\ref{prop:cm-small-primes}.
Theorem 2.8(2), p.~9, of \cite{Ishii} states that
\(\kappa_{\boldsymbol m}\) factors through \(A\).
Lemma 4.1, p.~17, gives its nonvanishing for every
\(\boldsymbol m\in I_K\).

Write \(\Lambda_2=\mathbb Z_p[[S,T]]\) for the group ring
of \(\operatorname{Gal}(L_\infty/K(p))\). The construction
in Lemma 4.3 and Proposition 4.4, pp.~19--20, gives an
equivariant surjection from one rank-one module for each
tame character onto its component of \(A\). The ordinary
source components are copies of \(\Lambda_2\); the exceptional
component is \(\operatorname{Ann}_{\Lambda_2}\mathbb Z_p(1)\).
Their total rank is \([K(p):K]\), which equals the rank of \(A\).
Since the source is torsion-free, the surjection has zero
kernel. Consequently every tame-character component of \(A\)
is torsion-free of rank one over \(\Lambda_2\).

We spell out the rank-one separation argument, including the
exceptional ideal. Fix a tame component \(M\), choose
\(0\ne v\in M\), and use its rank-one fraction space to choose
\(0\ne b\in\Lambda_2\) with
\[
 bM\subset\Lambda_2v .
\]
This is possible by clearing the denominators of finitely
many \(\Lambda_2\)-generators. For \(a\in M\), write
\(ba=f_av\), \(f_a\in\Lambda_2\).

Choose the two local generators so that their nontrivial
CM-character values are \(u=1+p\). For the fixed tame
component, the pairs
\[
 m_1\equiv r_1\pmod{p-1},\qquad
 m_2\equiv r_2\pmod{p-1}
\]
with both coordinates sufficiently large form a Cartesian
grid in \(I_K\). Set \(t_i=u^{m_i}-1\).
Equivariance gives
\[
 b(t_1,t_2)\kappa_{\boldsymbol m}(a)
       =f_a(t_1,t_2)\kappa_{\boldsymbol m}(v).
\]
At a grid point where \(b(t_1,t_2)\ne0\),
\(\kappa_{\boldsymbol m}(v)\ne0\): otherwise the same relation
for every generator of \(M\) would force the nonzero character
\(\kappa_{\boldsymbol m}\) to vanish identically.
Thus, when all characters kill \(a\), \(f_a\) vanishes at
every grid point at which \(b\) is nonzero. The product
\(bf_a\) vanishes at every grid point, including those where
\(b=0\).

A bounded one-variable power series has only finitely many
zeros in \(p\mathbb Z_p\), unless it is zero, by Weierstrass
preparation. Applying this fact first with \(t_2\) fixed,
then to the coefficients as \(t_2\) varies, proves that a
bounded two-variable series vanishing on this Cartesian
grid is zero. Hence \(bf_a=0\), so \(f_a=0\), and
torsion-freeness gives \(a=0\). Applying this argument to
each tame component proves the injectivity.

The elliptic Soul\'e characters are obtained from the
metabelian geometric outer action, so they kill the actual
outer kernel. Its image in \(A\) is therefore zero by the
injectivity just proved. This is the first containment.
The field containment follows by Galois correspondence.
\end{proof}

This proof uses the explicit module structure and the
nonzero characters. It does not use the freeness assertion
in Ishii's Conjecture 2.3 or the asserted converse in his
Proposition 2.11. The remaining subgroup
\(\overline{[H,H]}\) is explicit in the conclusion.

\subsection{The nonzero two-variable transgression}

\begin{proposition}\label{prop:cm-nonsplit}
For each of the two pairs in Proposition~\ref{prop:cm-small-primes},
put \(L=K(p)=K(\mu_p)\), \(L_\infty=\bigcup_nK(p^n)\),
\(G=\operatorname{Gal}(\Omega/L)\), and
\(H=\operatorname{Gal}(\Omega/L_\infty)\).
The exact sequence
\[
 1\longrightarrow H\longrightarrow G\longrightarrow
 \Gamma=\operatorname{Gal}(L_\infty/L)\simeq\mathbb Z_p^2
 \longrightarrow1
\]
has no continuous section. In fact, every pair of lifts to \(G\)
of a \(\mathbb Z_p\)-basis of \(\Gamma\) has a commutator detected
by a \(\Gamma\)-invariant homomorphism \(H\to\mathbb F_p\).
\end{proposition}
\begin{proof}
Let \(\Delta=\operatorname{Gal}(L/K)\), and write \(\omega\)
for its action on \(\mu_p\). Here \(\Delta\) has order \(p-1\)
and \(\omega\ne1\). Its action on \(\Gamma\) is trivial, since
\(L_\infty/K\) is abelian. A prime-to-\(p\) complement supplies
the \(\Delta\)-action on \(G\), and all invariant calculations
are independent of that choice.

Kummer theory and the already proved \(p\nmid h(L)\) give
\[
 H^1(G,\mu_p)\simeq\mathcal O_L[1/p]^\times/p .
\]
There are precisely two primes above \(p\), both fixed by
\(\Delta\). The rational \(S\)-unit representation is
\(\mathbb Q[\Delta]\oplus\mathbb Q\): the archimedean places
give the regular permutation representation, the two finite
places give two trivial representations, and the product
formula removes one trivial representation. Therefore the
nontrivial character \(\omega\) occurs once in its free part.
It occurs once more in the torsion subgroup \(\mu_p\).
Since \(p\nmid|\Delta|\), the integral character projectors
preserve these multiplicities upon reduction modulo \(p\).
Untwisting gives
\[
 \dim_{\mathbb F_p}H^1(G,\mathbb F_p)^\Delta=2.
\]

The comparison map from maximal pro-\(p\) group cohomology
to \'{e}tale cohomology is injective in degree two.
One can verify the needed injectivity by the five-term
sequence: a nonzero invariant homomorphism from the kernel
of the maximal pro-\(p\) quotient to a finite \(p\)-group
would produce a larger pro-\(p\) quotient. Thus
\[
 H^2(G,\mu_p)\hookrightarrow
 H^2_{\mathrm{\acute et}}(\mathcal O_L[1/p],\mu_p)
 \simeq\operatorname{Br}(\mathcal O_L[1/p])[p].
\]
The last isomorphism uses the vanishing of the class-group
term. Local Brauer invariants identify the last group with
\[
 \{(a,b)\in\mathbb F_p^2:a+b=0\}.
\]
The action of \(\Delta\) on these invariants is trivial,
because it fixes each of the two primes. On removing the
coefficient twist, it follows that
\[
 H^2(G,\mathbb F_p)^\Delta=0.
\]

The inflation map from \(H^1(\Gamma,\mathbb F_p)\) is
injective and both it and \(H^1(G,\mathbb F_p)^\Delta\)
have dimension two. It is consequently an isomorphism.
Taking \(\Delta\)-invariants in the Hochschild--Serre
five-term sequence now gives an isomorphism
\[
 d_2:\operatorname{Hom}(H,\mathbb F_p)^{\Gamma,\Delta}
       \xrightarrow{\ \sim\ }
       H^2(\Gamma,\mathbb F_p)\simeq\mathbb F_p .
\]
Choose \(\varphi\) with \(d_2\varphi\ne0\). The quotient
of \(G\) by \(\ker\varphi\) is the associated central
extension of \(\Gamma\) by \(\mathbb F_p\).
For coordinates \((a,b)\) on \(\Gamma\), the cocycle
\[
 ((a,b),(a',b'))\longmapsto ab'\pmod p
\]
represents a generator of \(H^2(\Gamma,\mathbb F_p)\).
Its alternating part has value one on a basis. Hence
\(\varphi([\widetilde\gamma_1,\widetilde\gamma_2])\ne0\)
for every pair of lifts of a basis. In particular the
lifts cannot commute, and a section cannot exist.
\end{proof}

This is the specific obstruction to using two commuting
arithmetic lifts in these examples. It agrees with Ishii's
construction, in which \([\gamma_1,\gamma_2]\) supplies the
trivial tame-character component of the abelianized kernel
\cite[\S4.1--4.2]{Ishii}. It is distinct from the
semilinear fixed-part calculation only in the explicit
sense established by the two displayed maps: the latter
calculates the coefficient-ring action, while \(d_2\)
calculates the extension class.

\subsection{The cyclotomic descent map}

Continue with either of the two explicit pairs, and put
\[
 L_\infty=K(p^\infty),\qquad
 L_{\mathrm c}=K(\mu_{p^\infty}),\qquad
 \Gamma_{\mathrm a}=\operatorname{Gal}(L_\infty/L_{\mathrm c}).
\]
The Weil-pairing identity \(\chi_1\chi_2=\chi_{\mathrm{cyc}}\)
identifies \(\Gamma_{\mathrm a}\) with the subgroup generated
by \(\gamma_1\gamma_2^{-1}\). Thus the actual group-ring quotient is
\[
 \mathbb Z_p[[S,T]]\longrightarrow\mathbb Z_p[[Z]],
 \qquad S\longmapsto Z,\quad T\longmapsto Z,
 \qquad \ker=(S-T).
\]
Complex conjugation exchanges \(\gamma_1,\gamma_2\), inverts
\(\Gamma_{\mathrm a}\), and acts trivially on the cyclotomic
quotient.

Put \(H_{\mathrm c}=\operatorname{Gal}(\Omega/L_{\mathrm c})\)
and \(\mathfrak X_{\mathrm c}=H_{\mathrm c}^{\mathrm{ab}}\).
The exact sequence
\(1\to H\to H_{\mathrm c}\to\Gamma_{\mathrm a}\to1\)
splits, since \(\Gamma_{\mathrm a}\simeq\mathbb Z_p\) is
projective in the category of pro-\(p\) groups. Abelianizing
this semidirect product gives the canonical exact sequence
\[
 0\longrightarrow A_{\Gamma_{\mathrm a}}
  \longrightarrow\mathfrak X_{\mathrm c}
  \longrightarrow\Gamma_{\mathrm a}\longrightarrow0.
\]
Taking fixed parts for complex conjugation is exact, because
two is invertible, and \(\Gamma_{\mathrm a}^+=0\). Hence
\[
 \mathfrak X_{\mathrm c}^+
 \simeq\bigl(A/(S-T)A\bigr)^+
 \simeq A^+/(S-T)A^- .
\]
For the last equality, every fixed quotient class has a
fixed lift by averaging. Moreover, if \((S-T)a\) is fixed,
it equals \((S-T)(a-ca)/2\). This proves the asserted kernel
without assuming that \(A^+\) is a module over the whole
two-variable ring.

Let \(X_{\mathrm c}\) be the unramified abelian pro-\(p\)
module over \(L_{\mathrm c}\), and let \(I_{\mathrm c}\)
be the closed subgroup of \(\mathfrak X_{\mathrm c}\)
generated by inertia at the primes above \(p\).
The exact ramification quotient is
\[
 0\longrightarrow I_{\mathrm c}^+
 \longrightarrow\mathfrak X_{\mathrm c}^+
 \longrightarrow X_{\mathrm c}^+\longrightarrow0.
\]
Thus the class-group fixed part is reached through both
the displayed anti-cyclotomic coinvariants and this
explicit inertia quotient.

\begin{proposition}\label{prop:cm-unramified-zero}
For the two explicit examples, \(X_{\mathrm c}^+=0\).
In addition, the full unramified abelian pro-\(p\)
module over \(L_\infty\) is zero.
\end{proposition}
\begin{proof}
Write \(L_r=K(\mu_{p^r})\) and \(L_r^+\) for its real
subfield. Extension of ideals and norm identify the
\(p\)-primary class group of \(L_r^+\) with the fixed
part of that of \(L_r\): their composites are multiplication
by two, and extension followed by norm on a fixed class
is also multiplication by two. The initial real class
group therefore has order prime to \(p\).

Each extension \(L_{r+1}^+/L_r^+\) is cyclic of degree
\(p\) and has exactly one ramified prime, the unique
prime above \(p\), which is totally ramified.
Indeed \(p\) splits into two primes in \(K\), and complex
conjugation exchanges the two primes in \(L_r\);
passing to the real field identifies them. The local
cyclotomic tower is totally ramified.

Here is the class-field induction. In a cyclic degree-\(p\)
extension \(F'/F\) with just one totally ramified prime,
suppose the \(p\)-class group of \(F\) is zero.
The Hilbert \(p\)-class field \(B\) of \(F'\) is Galois
over \(F\), with a finite \(p\)-group \(P\). Its inertia
subgroup at that prime has order \(p\), since \(B/F'\)
is unramified. The normal closure of this inertia is
all of \(P\): a nontrivial unramified quotient would
have a cyclic quotient of order \(p\), contrary to
the assumption on \(F\). Thus \(P\) is normally
generated by one element, so \(P/\Phi(P)\) is cyclic,
and the Burnside basis theorem makes \(P\) cyclic.
Its inertia is then \(P\), which has order \(p\);
therefore \(B=F'\). Induction proves that every
\(p\)-class group of \(L_r^+\) is zero. Passing to
norm limits gives \(X_{\mathrm c}^+=0\).

For the second assertion, use the proved locality
isomorphism of Proposition~\ref{prop:cm-locality}.
Let \(I_{\mathfrak p}\subset G\) be local inertia.
Then \(G/I_{\mathfrak p}\simeq\mathbb Z_p\), the
maximal pro-\(p\) unramified local quotient.
Under \(G\to\Gamma=\operatorname{Gal}(L_\infty/L)\),
the inertia image is exactly the first principal-unit
factor \(\mathbb Z_p\times0\). This is the ray-class
reciprocity map; the quotient is the second
\(\mathbb Z_p\)-factor. Since decomposition is all
of \(\Gamma\), the induced map
\[
 G/I_{\mathfrak p}\longrightarrow
 \Gamma/\operatorname{im}(I_{\mathfrak p})
\]
is a surjection between two copies of \(\mathbb Z_p\),
and hence an isomorphism. Its kernel calculation
gives \(H\subset I_{\mathfrak p}\).

Every finite unramified pro-\(p\) extension of
\(L_\infty\) descends to a finite layer as an extension
unramified outside \(p\); its normal closure over \(L\)
is still a \(p\)-extension and lies in \(\Omega\).
In \(\Omega/L_\infty\), the inertia subgroup at
\(\mathfrak p\) is \(H\cap I_{\mathfrak p}=H\).
An unramified quotient must kill it and is therefore
trivial. This proves the assertion about the full
unramified abelian module.
\end{proof}

\subsection{The two-variable cup edge}

The extension class affects the precise cohomological map.
For a finite coefficient module \(B\) on which \(H\) acts
trivially, the Hochschild--Serre sequence for
\(1\to H\to G\to\Gamma\to1\), using
\(\operatorname{cd}_p(H)=1\) and
\(\operatorname{cd}_p(\Gamma)=2\), gives
\[
 0\longrightarrow
 \operatorname{coker}\!\left(
  H^1(H,B)^\Gamma\xrightarrow{d_2}H^2(\Gamma,B)\right)
 \longrightarrow H^2(G,B)
 \longrightarrow H^1(\Gamma,H^1(H,B))
 \longrightarrow0 .
\]
The first arrow is inflation. The differential \(d_2\)
is the extension transgression calculated in
Proposition~\ref{prop:cm-nonsplit}.
For a \(\Gamma\)-module \(M\), put \(D_i=\gamma_i-1\).
The two-variable Koszul complex computes the right
side as
\[
 H^1(\Gamma,M)=
 \frac{\{(m_1,m_2):D_2m_1=D_1m_2\}}
      {\{(D_1m,D_2m):m\in M\}} .
\]
For \(B=\mathbb Z/p^n(2)\), complex conjugation acts on
this complex by \(c\) in degree zero, by
\((m_1,m_2)\mapsto(cm_2,cm_1)\) in degree one,
and by \(-c\) in degree two. The sign in degree two
comes from exchanging the two generators.
In particular the fixed part of the inflated term is
zero, whereas the fixed part of the edge is represented
by the equation
\[
 D_2m=D_1cm
\]
modulo \(D_1M^+\). These explicit maps identify what
must replace the procyclic coinvariant calculation
in a two-variable cup argument.

\subsection{An explicit nonzero commutator in finite geometric quotients}
\label{subsec:cm-low-class-two}

This computation specifies a small nonabelian quotient of the
outer image.  Its definition does not assert detection of the full
class-two quotient of the arithmetic source.
Let \(\Pi=\pi_1^{(p)}((E-\{O\})_{\bar K})\), with descending central
series \(\Pi(n)\), and put
\[
 P=\Pi/\Pi(13),\qquad
 P_r=P/\overline{\langle x^{p^r}:x\in P\rangle}\quad(r\geq1).
\]
These are characteristic quotients of \(\Pi\). The successive
lower-central quotients of \(P_r\) are finitely generated abelian
groups of exponent dividing \(p^r\), so every \(P_r\) is a finite \(p\)-group.
The field identity~\eqref{eq:cm-radical-outer-field} gives the
actual induced outer map \(\rho_H:H\longrightarrow\operatorname{Out}(\Pi)\):
indeed, restriction identifies
\(\operatorname{Gal}(K(E[p])\Omega/K(E[p]))\) with
\(\operatorname{Gal}(\Omega/L_1)\), because the intersection is
\(L_1\), and the outer action factors through its quotient
\(\operatorname{Gal}(\Omega^*/L_1)\).
Define the finite image
\[
 Q_r=\operatorname{im}\bigl(H\xrightarrow{\rho_H}
                              \operatorname{Out}(P_r)\bigr).
\]

\begin{proposition}\label{prop:cm-small-nonabelian-image}
For every class-number-one imaginary quadratic \(K\), every
split \(p\geq5\), and every \(E/K\) as above, all \(Q_r\) have
nilpotency class at most two.  They are nonabelian for all
sufficiently large \(r\).
\end{proposition}

\begin{proof}
We first give the rational Lie calculation.
Let \(L\) be the free Lie algebra over \(\mathbb Q\) on \(a,b\).
Let \(D_{rs}\) denote its space of derivations with
\[
 D(a)\in L_{r+1,s},\qquad D(b)\in L_{r,s+1},\qquad
                         D([a,b])=0.
\]
For a Lyndon word \(w\), write \([w]\) for its standard bracket:
split off its longest proper Lyndon suffix and bracket
recursively, with \(a<b\).
The following coefficient vectors define \(d\in D_{22}\) and
\(e,f\in D_{33}\).  The entries of each row are coefficients in
the displayed ordered Lyndon basis:
\[
\begin{array}{c|c|r@{\,,\,}r@{\,,\,}r@{\,,\,}r@{\,,\,}r}
 &\text{ordered basis}&\multicolumn{5}{c}{\text{coefficients}}\\ \hline
d(a)&[aaabb],[aabab]&\multicolumn{5}{c}{(-2,-1)}\\
d(b)&[aabbb],[ababb]&\multicolumn{5}{c}{(2,3)}\\
e(a)&[aaaabbb],[aaababb],[aaabbab],[aabaabb],[aababab]
       &1&-5&-6&1&-2\\
e(b)&[aaabbbb],[aababbb],[aabbabb],[aabbbab],[abababb]
       &-1&4&6&8&0\\
f(a)&[aaaabbb],[aaababb],[aaabbab],[aabaabb],[aababab]
       &-3&-9&-6&-3&-2\\
f(b)&[aaabbbb],[aababbb],[aabbabb],[aabbbab],[abababb]
       &3&12&6&0&8
\end{array}
\]
Direct expansion gives
\(d([a,b])=e([a,b])=f([a,b])=0\).
The boundary map
\[
 L_{r+1,s}\oplus L_{r,s+1}\longrightarrow L_{r+1,s+1},
 \qquad (u,v)\longmapsto[u,b]+[a,v],
\]
is surjective: a free Lie algebra generated by \(a,b\) is the
sum of \([a,L]\) and \([b,L]\) in every degree greater than one.
The multigraded Witt formula gives
\[
 \dim L_{3,2}=\dim L_{2,3}=2,\quad \dim L_{3,3}=3,
 \qquad
 \dim L_{4,3}=\dim L_{3,4}=5,\quad \dim L_{4,4}=8.
\]
Thus \(\dim D_{22}=1\) and \(\dim D_{33}=2\).
The displayed nonzero \(d\) spans \(D_{22}\);
the distinct last two \(b\)-coefficient vectors show that \(e,f\)
are independent and hence form a basis of \(D_{33}\).

To verify injectivity of the bracket map, only two associative
word coefficients are needed.  Expansion by the derivation rule
gives
\[
\begin{array}{c|rr}
 \text{coefficient in the image of }a &[d,e]&[d,f]\\ \hline
 aaaaabbabbb&-2&6\\
 aaaababbabb&-22&18
\end{array}
\qquad
\det\begin{pmatrix}-2&6\\-22&18\end{pmatrix}=96.
\]
Consequently \([d,\cdot]:D_{33}\longrightarrow D_{55}\) is
injective over \(\mathbb Q_p\), for every \(p\).
All equalities here are integral free-associative-algebra
identities.  The accompanying exact program
\texttt{check\_cm\_low\_class\_two.py} reconstructs the Lyndon
bases, checks the boundary equations, and records the complete
coefficient certificate.

We now identify the arithmetic classes to which this calculation
applies.  By \cite[Theorem 2.8 and Lemma 2.10, p.~9]{Ishii},
\(\kappa_{(2,2)}\) and \(\kappa_{(3,3)}\) are nonzero on the
weight-four and weight-six graded quotients, respectively.
Their nonvanishing uses the proven finiteness of the
second cohomology groups with Tate twists two and three.
The CM character decomposition of the rational graded pieces
is semisimple by \cite[Proposition 2.2]{Ishii}.
It therefore supplies nonzero rational classes of pure
bidegrees \((2,2)\) and \((3,3)\).
Clearing their denominators gives integral graded classes,
represented by actual elements \(h_4\in F^4H\) and
\(h_6\in F^6H\).
These elements are not asserted to satisfy exact eigen-equations
under conjugation.

The symplectic derivation comparison
\cite[Proposition 3.3(4), pp.~12--14]{IshiiLie} identifies these
two outer graded classes with nonzero elements of \(D_{22}\)
and \(D_{33}\).  Its inner kernel is generated by
\(\operatorname{ad}[a,b]\) in bidegree \((1,1)\);
it has zero component in bidegrees \((2,2),(3,3),(5,5)\).
The calculation therefore proves that
\([h_4,h_6]\) has nonzero outer class in weight ten.
In particular its image in \(\operatorname{Out}(P)\) is nontrivial.

By \cite[Proposition 2.2(1)]{Ishii}, the weight-one,
weight-two, and weight-three arithmetic graded pieces vanish.
Thus \(\rho_H(H)\subseteq F^4\operatorname{Out}(\Pi)\).
Centrality of the weight filtration gives
\[
 \rho_H(\gamma_3H)\subseteq F^{12}\operatorname{Out}(\Pi)
       =\ker\bigl(\operatorname{Out}(\Pi)\longrightarrow
                              \operatorname{Out}(\Pi/\Pi(13))\bigr).
\]
Every \(Q_r\) consequently has class at most two.

Finally, the natural map
\(\operatorname{Out}(P)\longrightarrow
 \varprojlim_r\operatorname{Out}(P_r)\) is injective.
To see this, suppose an automorphism is inner modulo every
power subgroup defining \(P_r\).  For each \(r\), its possible
conjugating elements form a nonempty closed subset of the
compact group \(P\), and these subsets are nested.
Their intersection contains a conjugator on \(P\), since
the intersection of those power subgroups is trivial.
The latter assertion follows because every finite quotient of
the pro-\(p\) group \(P\) has finite \(p\)-power exponent.
Hence the nontrivial outer image of \([h_4,h_6]\) survives
in some \(\operatorname{Out}(P_r)\) and in every later one.
This proves the proposition.
\end{proof}

The proposition gives an explicit nonzero central bracket in
the actual finite outer images.  Its source is the displayed
homomorphism \(H\to Q_r\); it does not identify the kernel of
the different homomorphism \(H\to H/\gamma_3H\).


\subsection{What the geometric common kernel gives}
\label{subsec:cm-geometric-common-kernel}

Write \(\Pi=\widehat{\langle x,y\rangle}\) for the geometric
fundamental group of a once-punctured elliptic curve and
\[
 R_{\mathrm{geom}}=
 \bigcap_\ell\ker(\Pi\longrightarrow\Pi^{(\ell)}).
\]
For cusp inertia \(I=\overline{\langle[x,y]\rangle}\), one has
\(I\cap R_{\mathrm{geom}}=1\). Indeed, for every \(n\geq1\)
the homomorphism to \(\operatorname{UT}_3(\mathbb Z/\ell^n)\)
sending \(x\) to \(1+E_{12}\) and \(y\) to \(1+E_{23}\)
sends \([x,y]\) to \(1+E_{13}\), of order \(\ell^n\).
The map \(\widehat{\mathbb Z}\to I\) is consequently injective:
each nonzero component \(\alpha_\ell\) is detected modulo some
\(\ell^n\). This also proves the asserted intersection, for every
conjugate of \(I\).

The arithmetic outer kernel occurs in the quotient row of the
fundamental exact sequence. More precisely, for
\[
 1\longrightarrow\Pi^{(p)}\longrightarrow
 \pi_1(X)^{(p,\mathrm{geom})}\longrightarrow G_K\longrightarrow1
\]
and \(N=\ker(\rho_{X,p})\), the inverse image of \(N\) equals
\(\Pi^{(p)}C\), where \(C\) is the centralizer of \(\Pi^{(p)}\)
in the middle group. To check this equality, take an element
above \(N\). Its action on \(\Pi^{(p)}\) is conjugation by some
\(u\in\Pi^{(p)}\); multiplication by \(u^{-1}\) places it in \(C\).
The converse follows by reversing this operation.
Since the center of the rank-two free pro-\(p\) group is trivial,
\(C\to N\) is an isomorphism. Thus Theorem~\ref{thm:geometric-procyclic-recognition} identifies the
procyclic subgroups of \(R_{\mathrm{geom}}\) with open arithmetic
normalizer image, while the displayed isomorphism places \(N\)
in the centralizer \(C\). A proof using that
recognition assertion to annihilate \(N\) still requires a map
from a nontrivial element of \(C\) to a subgroup to which the
recognition assertion applies. No such map is established here.

The limitation of that particular input also has an explicit
group-theoretic test. Start with an action \(G\to\operatorname{Aut}(\Pi)\)
and its extension \(\Pi\rtimes G\), and choose a nonzero pro-\(\ell\)
group \(D\). Pull back by the projection
\(q:G\times D\to G\). The pulled-back extension is
\((\Pi\rtimes G)\times D\). For every subgroup \(A\subset\Pi\),
its normalizer in this extension is
\(N_{\Pi\rtimes G}(A)\times D\). Its image in \(G\times D\)
is therefore the inverse image under \(q\) of the old normalizer
image. Such an inverse image is open exactly when the original
image is open: one implication is continuity, and the converse
follows because \(q\) is an open quotient map. Every recognition
assertion expressed by these normalizer images is consequently
preserved. The new outer action is \(\rho\circ q\), with kernel
\(\ker(\rho)\times D\). Thus this recognition input alone leaves
an arbitrary factor \(D\) in the arithmetic outer kernel.
This calculation concerns the stated group extension; it claims
no realization of the pullback as a number-field Galois action.

\subsection{The remaining elliptic gate}

For the two examples just computed, Ishii's generation theorem
applies: its cohomological input follows from the two arithmetic
hypotheses by \cite[Lemma 1.5, p.~3]{IshiiLie}.
The source group is free by Proposition~\ref{prop:cm-finite-plus-free}
and the vanishing proved above for these examples. Freeness of the
arithmetic graded Lie algebra in Ishii's proposed presentation remains
open for these two fields.
For the other class-number-one CM fields, the explicit Pollack
relation above proves that the proposed free presentation must
be replaced by a presentation incorporating relations.
The weighted-completion map is a surjection onto the unipotent
outer image; its injectivity is precisely an additional assertion,
as the known Pollack relation demonstrates in the indicated
bidegree. This gives a specific morphism to test, rather than
an inference from the existence of elliptic polylogarithms.

Goto and Ishii's more recent Coleman--Ihara formula
\cite[Theorem 1.1 and Corollary 5.5]{GotoIshii2026}
expresses restrictions of elliptic Soul\'e characters to local
reciprocity images through two-variable Coates--Wiles
homomorphisms and Katz \(p\)-adic \(L\)-values. Its domain is the
norm limit of semilocal principal units, and its target is a
one-dimensional character. Sano and Venjakob's extended
interpolation formula
\cite[Theorem 1.2.1, p.~3, and Theorem 2.4.1, p.~8]{SanoVenjakob2026}
uses \(V_\chi=\chi_{\mathrm{cyc}}\chi^{-1}\) with
\(\chi=\psi^k\bar\psi^j\). The Weil-pairing identity
\(\psi\bar\psi=\chi_{\mathrm{cyc}}\) identifies the exponents of
\(V_\chi\) with \((1-k,1-j)\). Hence the positive target
\((m_1,m_2)\) has \(k=1-m_1\leq0\), \(j=1-m_2\leq0\).
Its comparison is the noncritical one in
\cite[Corollary 2.4.7, pp.~9--10]{SanoVenjakob2026};
the source does not assert nonvanishing of every such regulator.
Theorem~\ref{thm:cm-finite-plus-abelian} gives the abelian component
fields through the exact Kummer pairing. Full nonabelian faithfulness
remains open, including for the two examples. A homogeneous Pollack
relation gives an element whose outer image lies deeper in the
weight filtration. An element of the group kernel requires a correction
inside the correspondingly deeper filtration of the source. The required
strictness of the filtered outer map is not proved here: rational
generation of each graded piece allows denominators that can grow
without bound during successive corrections.
